\input{preamble}

% OK, start here.
%
\begin{document}

\title{Espaces algébriques}


\maketitle

\phantomsection
\label{section-phantom}

\tableofcontents

\section{Introduction}
\label{section-introduction}

\noindent
Les espaces algébriques ont été introduits par Michael Artin,
voir \cite{ArtinI}, \cite{ArtinII},
\cite{Artin-Theorem-Representability},
\cite{Artin-Construction-Techniques},
\cite{Artin-Algebraic-Spaces},
\cite{Artin-Algebraic-Approximation},
\cite{Artin-Implicit-Function},
et \cite{ArtinVersal}.
Une partie des fondements a été élaborée conjointement avec
Knutson, qui en a tiré l'ouvrage \cite{Kn}.
Artin a défini (voir \cite[Définition 1.3]{Artin-Implicit-Function})
un espace algébrique comme un faisceau pour la topologie étale
qui est localement représentable pour la topologie étale.
Dans la plupart des travaux d'Artin, les catégories de schémas
considérées sont formées de schémas localement de type fini sur une base
noethérienne excellente fixée.

\medskip\noindent
Notre définition diffère légèrement de la définition originelle d'Artin.
En effet, nos espaces algébriques sont des faisceaux pour la topologie fppf
dont la diagonale est représentable et qui possèdent un «~revêtement~» étale
par un schéma. Travailler avec la topologie fppf plutôt qu'avec la topologie
étale n'est qu'un point technique et ne change presque
rien; nous montrerons dans
Amorçage, section \ref{bootstrap-section-spaces-etale}
que nous aurions obtenu la même catégorie d'espaces algébriques
en travaillant avec la topologie étale. Dans ce même chapitre,
nous démontrerons que la condition portant sur la diagonale
peut, en un certain sens, être supprimée; voir
Amorçage, section \ref{bootstrap-section-bootstrap}.

\medskip\noindent
Après avoir défini les espaces algébriques, nous formulons quelques
observations fondamentales. Le principal résultat de ce chapitre est qu'avec
nos définitions, la donnée d'un espace algébrique équivaut à celle d'une
relation d'équivalence étale; voir l'exposé de la section
\ref{section-presentations} et le Théorème \ref{theorem-presentation}.
L'analogue de ce théorème dans le cadre d'Artin est \cite[Théorème 1.5]{Artin-Implicit-Function}, ou
\cite[Proposition II.1.7]{Kn}. Autrement dit, le faisceau
défini par une relation d'équivalence étale a une diagonale représentable.
Il s'ensuit que notre définition coïncide, au sens large, avec la définition
originelle d'Artin. Cela signifie aussi que l'on peut donner des exemples
d'espaces algébriques en écrivant simplement une relation d'équivalence étale.

\medskip\noindent
Dans la section \ref{section-separation}, nous introduisons divers axiomes
de séparation pour les espaces algébriques rencontrés dans la littérature.
Enfin, dans la section \ref{section-examples},
nous donnons des exemples insolites d'espaces algébriques et d'autres qui le sont moins.






\section{Remarques générales}
\label{section-general}

\noindent
Nous travaillons dans un grand site fppf convenable $\Sch_{fppf}$
comme dans Topologies, Définition \ref{topologies-definition-big-fppf-site}.
Ainsi, sauf mention expresse du contraire, tous les schémas seront des objets
de $\Sch_{fppf}$.
Dans la section \ref{section-change-big-site}, nous étudions les modifications
entraînées par un changement de grand site fppf.

\medskip\noindent
Nous travaillerons toujours relativement à un schéma de base $S$ appartenant à $\Sch_{fppf}$.
Nous travaillerons alors avec le grand site fppf $(\Sch/S)_{fppf}$,
voir Topologies, Définition \ref{topologies-definition-big-small-fppf}.
On retrouve le cas absolu en prenant
$S = \Spec(\mathbf{Z})$.

\medskip\noindent
Si $U, T$ sont des schémas sur $S$, nous désignons
par $U(T)$ l'ensemble des points à valeurs dans $T$ {\it au-dessus} de $S$.
Autrement dit : $U(T) = \Mor_S(T, U)$.

\medskip\noindent
Notons que tout recouvrement fpqc est un épimorphisme effectif universel, voir
Descente, Lemme \ref{descent-lemma-fpqc-universal-effective-epimorphisms}.
Par conséquent, la topologie sur $\Sch_{fppf}$
est moins fine que la topologie canonique, et tous les préfaisceaux représentables
sont des faisceaux.







\section{Morphismes représentables de préfaisceaux}
\label{section-representable}

\noindent
Soit $S$ un schéma appartenant à $\Sch_{fppf}$.
Soient $F, G : (\Sch/S)_{fppf}^{opp} \to \textit{Ens}$.
Soit $a : F \to G$ un morphisme représentable de foncteurs, voir
Catégories,
Définition \ref{categories-definition-representable-map-presheaves}.
Cela signifie que pour tout
$U \in \Ob((\Sch/S)_{fppf})$
et tout $\xi \in G(U)$, le produit fibré $h_U \times_{\xi, G} F$ est représentable.
Choisissons un objet $V_\xi$ qui le représente et un isomorphisme
$h_{V_\xi} \to h_U \times_G F$.
D'après le lemme de Yoneda, voir Catégories, Lemme \ref{categories-lemma-yoneda},
la projection $h_{V_\xi} \to h_U \times_G F \to h_U$ provient d'un unique
morphisme de schémas $a_\xi : V_\xi \to U$.
Le diagramme suivant permet de se représenter la situation :
$$
\xymatrix{
V_\xi \ar@{~>}[r] \ar[d]_{a_\xi} & h_{V_\xi} \ar[d] \ar[r] & F \ar[d]^a \\
U \ar@{~>}[r] & h_U \ar[r]^\xi & G
}
$$
où les flèches ondulées représentent le plongement de Yoneda.
Voici quelques lemmes sur cette notion, valables dans un cadre très général.

\begin{lemma}
\label{lemma-morphism-schemes-gives-representable-transformation}
Soit $S$ un schéma appartenant à $\Sch_{fppf}$ et soient
$X$, $Y$ des objets de $(\Sch/S)_{fppf}$.
Soit $f : X \to Y$ un morphisme de schémas.
Alors
$$
h_f : h_X \longrightarrow h_Y
$$
est un morphisme représentable de foncteurs.
\end{lemma}

\begin{proof}
Cela est formel et repose seulement sur le fait que
la catégorie $(\Sch/S)_{fppf}$ admet des produits fibrés.
\end{proof}

\begin{lemma}
\label{lemma-composition-representable-transformations}
Soit $S$ un schéma appartenant à $\Sch_{fppf}$.
Soient $F, G, H : (\Sch/S)_{fppf}^{opp} \to \textit{Ens}$.
Soient $a : F \to G$, $b : G \to H$ des morphismes représentables de foncteurs.
Alors
$$
b \circ a : F \longrightarrow H
$$
est un morphisme représentable de foncteurs.
\end{lemma}

\begin{proof}
Cela est entièrement formel et vaut dans toute catégorie.
\end{proof}

\begin{lemma}
\label{lemma-base-change-representable-transformations}
Soit $S$ un schéma appartenant à $\Sch_{fppf}$.
Soient $F, G, H : (\Sch/S)_{fppf}^{opp} \to \textit{Ens}$.
Soit $a : F \to G$ un morphisme représentable de foncteurs.
Soit $b : H \to G$ un morphisme de foncteurs quelconque.
Considérons le diagramme de produit fibré
$$
\xymatrix{
H \times_{b, G, a} F \ar[r]_-{b'} \ar[d]_{a'} & F \ar[d]^a \\
H \ar[r]^b & G
}
$$
Alors le changement de base $a'$ est un morphisme représentable de foncteurs.
\end{lemma}

\begin{proof}
Cela est entièrement formel et vaut dans toute catégorie.
\end{proof}

\begin{lemma}
\label{lemma-product-representable-transformations}
Soit $S$ un schéma appartenant à $\Sch_{fppf}$.
Soient $F_i, G_i : (\Sch/S)_{fppf}^{opp} \to \textit{Ens}$, $i = 1, 2$.
Soient $a_i : F_i \to G_i$, $i = 1, 2$,
des morphismes représentables de foncteurs.
Alors
$$
a_1 \times a_2 : F_1 \times F_2 \longrightarrow G_1 \times G_2
$$
est un morphisme représentable de foncteurs.
\end{lemma}

\begin{proof}
Écrivons $a_1 \times a_2$ comme la composée
$F_1 \times F_2 \to G_1 \times F_2 \to G_1 \times G_2$.
La première flèche est le changement de base de $a_1$ par le morphisme
$G_1 \times F_2 \to G_1$, et la seconde flèche
est le changement de base de $a_2$ par le morphisme
$G_1 \times G_2 \to G_2$. Ce lemme est donc une conséquence
formelle des Lemmes \ref{lemma-composition-representable-transformations}
et \ref{lemma-base-change-representable-transformations}.
\end{proof}

\begin{lemma}
\label{lemma-representable-transformation-to-sheaf}
Soit $S$ un schéma appartenant à $\Sch_{fppf}$.
Soient $F, G : (\Sch/S)_{fppf}^{opp} \to \textit{Ens}$.
Soit $a : F \to G$ un morphisme représentable de foncteurs.
Si $G$ est un faisceau, alors $F$ l'est aussi.
\end{lemma}

\begin{proof}
Soit $\{\varphi_i : T_i \to T\}$ une famille couvrante du site
$(\Sch/S)_{fppf}$.
Soient $s_i \in F(T_i)$ satisfaisant à la condition de recollement.
Alors $\sigma_i = a(s_i) \in G(T_i)$ satisfont elles aussi à la condition de recollement.
Il existe donc un unique $\sigma \in G(T)$ tel
que $\sigma_i = \sigma|_{T_i}$. Par hypothèse,
$F' = h_T \times_{\sigma, G, a} F$ est un préfaisceau représentable
et donc (voir les remarques de la section \ref{section-general}) un faisceau.
Notons que $(\varphi_i, s_i) \in F'(T_i)$ satisfont
elles aussi à la condition de recollement et proviennent donc d'un unique
$(\text{id}_T, s) \in F'(T)$. Manifestement, $s$ est la section de
$F$ que nous cherchions.
\end{proof}

\begin{lemma}
\label{lemma-representable-transformation-diagonal}
Soit $S$ un schéma appartenant à $\Sch_{fppf}$.
Soient $F, G : (\Sch/S)_{fppf}^{opp} \to \textit{Ens}$.
Soit $a : F \to G$ un morphisme représentable de foncteurs.
Alors $\Delta_{F/G} : F \to F \times_G F$ est représentable.
\end{lemma}

\begin{proof}
Soit $U \in \Ob((\Sch/S)_{fppf})$. Soit
$\xi = (\xi_1, \xi_2) \in (F \times_G F)(U)$.
Posons $\xi' = a(\xi_1) = a(\xi_2) \in G(U)$.
Par hypothèse, il existe un schéma $V$ muni d'un morphisme $V \to U$
qui représente le produit fibré $h_U \times_{\xi', G} F$.
En particulier, les éléments $\xi_1, \xi_2$ définissent des morphismes
$f_1, f_2 : U \to V$ au-dessus de $U$. Puisque $V$ représente le
produit fibré $h_U \times_{\xi', G} F$ et que
$\xi' = a \circ \xi_1 = a \circ \xi_2$
nous voyons que, si $g : U' \to U$ est un morphisme, alors
$$
g^*\xi_1 = g^*\xi_2
\Leftrightarrow
f_1 \circ g = f_2 \circ g.
$$
Autrement dit, nous voyons que $h_U \times_{\xi, F \times_G F} F$
est représenté par $V \times_{\Delta, V \times V, (f_1, f_2)} U$
qui est un schéma.
\end{proof}










\section{Listes de propriétés utiles des morphismes de schémas}
\label{section-lists}

\noindent
Pour faciliter les renvois, les remarques suivantes énumèrent les
propriétés de morphismes qui vérifient certaines des conditions
requises dans des résultats ultérieurs.

\begin{remark}
\label{remark-list-properties-stable-base-change}
Voici une liste de propriétés/types de morphismes
qui sont {\it stables par changement de base quelconque} :
\begin{enumerate}
\item immersions fermées, ouvertes et localement fermées, voir
Schémas, Lemme \ref{schemes-lemma-base-change-immersion},
\item quasi-compact, voir
Schémas, Lemme \ref{schemes-lemma-quasi-compact-preserved-base-change},
\item universellement fermé, voir
Schémas, Définition \ref{schemes-definition-universally-closed},
\item (quasi-)séparé, voir
Schémas, Lemme \ref{schemes-lemma-separated-permanence},
\item monomorphisme, voir
Schémas, Lemme \ref{schemes-lemma-base-change-monomorphism}
\item surjectif, voir
Morphismes, Lemme \ref{morphisms-lemma-base-change-surjective},
\item radiciel, voir
Morphismes, Lemme \ref{morphisms-lemma-universally-injective},
\item affine, voir
Morphismes, Lemme \ref{morphisms-lemma-base-change-affine},
\item quasi-affine, voir
Morphismes, Lemme \ref{morphisms-lemma-base-change-quasi-affine},
\item (localement) de type fini, voir
Morphismes, Lemme \ref{morphisms-lemma-base-change-finite-type},
\item (localement) quasi-fini, voir
Morphismes, Lemme \ref{morphisms-lemma-base-change-quasi-finite},
\item (localement) de présentation finie, voir
Morphismes, Lemme \ref{morphisms-lemma-base-change-finite-presentation},
\item localement de type fini de dimension relative $d$, voir
Morphismes, Lemme \ref{morphisms-lemma-base-change-relative-dimension-d},
\item universellement ouvert, voir
Morphismes, Définition \ref{morphisms-definition-open},
\item plat, voir
Morphismes, Lemme \ref{morphisms-lemma-base-change-flat},
\item syntomique, voir
Morphismes, Lemme \ref{morphisms-lemma-base-change-syntomic},
\item lisse, voir
Morphismes, Lemme \ref{morphisms-lemma-base-change-smooth},
\item non ramifié (resp.\ G-non ramifié), voir
Morphismes, Lemme \ref{morphisms-lemma-base-change-unramified},
\item \'etale, voir
Morphismes, Lemme \ref{morphisms-lemma-base-change-etale},
\item propre, voir
Morphismes, Lemme \ref{morphisms-lemma-base-change-proper},
\item H-projectif, voir
Morphismes, Lemme \ref{morphisms-lemma-H-projective-base-change},
\item (localement) projectif, voir
Morphismes, Lemme \ref{morphisms-lemma-base-change-projective},
\item fini ou entier, voir
Morphismes, Lemme \ref{morphisms-lemma-base-change-finite},
\item fini localement libre, voir
Morphismes, Lemme \ref{morphisms-lemma-base-change-finite-locally-free},
\item universellement submersif, voir
Morphismes, Lemme \ref{morphisms-lemma-base-change-universally-submersive},
\item homéomorphisme universel, voir
Morphismes, Lemme \ref{morphisms-lemma-base-change-universal-homeomorphism}.
\end{enumerate}
Compléter au besoin.
\end{remark}

\begin{remark}
\label{remark-list-properties-stable-composition}
Parmi les propriétés de morphismes qui sont stables par changement de base
(telles qu'elles sont énumérées dans la
Remarque \ref{remark-list-properties-stable-base-change}),
les suivantes sont aussi {\it stables par composition} :
\begin{enumerate}
\item immersions fermées, ouvertes et localement fermées, voir
Schémas, Lemme \ref{schemes-lemma-composition-immersion},
\item quasi-compact, voir
Schémas, Lemme \ref{schemes-lemma-composition-quasi-compact},
\item universellement fermé, voir
Morphismes, Lemme \ref{morphisms-lemma-composition-proper},
\item (quasi-)séparé, voir
Schémas, Lemme \ref{schemes-lemma-separated-permanence},
\item monomorphisme, voir
Schémas, Lemme \ref{schemes-lemma-composition-monomorphism},
\item surjectif, voir
Morphismes, Lemme \ref{morphisms-lemma-composition-surjective},
\item radiciel, voir
Morphismes, Lemme \ref{morphisms-lemma-composition-universally-injective},
\item affine, voir
Morphismes, Lemme \ref{morphisms-lemma-composition-affine},
\item quasi-affine, voir
Morphismes, Lemme \ref{morphisms-lemma-composition-quasi-affine},
\item (localement) de type fini, voir
Morphismes, Lemme \ref{morphisms-lemma-composition-finite-type},
\item (localement) quasi-fini, voir
Morphismes, Lemme \ref{morphisms-lemma-composition-quasi-finite},
\item (localement) de présentation finie, voir
Morphismes, Lemme \ref{morphisms-lemma-composition-finite-presentation},
\item universellement ouvert, voir
Morphismes, Lemme \ref{morphisms-lemma-composition-open},
\item plat, voir
Morphismes, Lemme \ref{morphisms-lemma-composition-flat},
\item syntomique, voir
Morphismes, Lemme \ref{morphisms-lemma-composition-syntomic},
\item lisse, voir
Morphismes, Lemme \ref{morphisms-lemma-composition-smooth},
\item non ramifié (resp.\ G-non ramifié), voir
Morphismes, Lemme \ref{morphisms-lemma-composition-unramified},
\item \'etale, voir
Morphismes, Lemme \ref{morphisms-lemma-composition-etale},
\item propre, voir
Morphismes, Lemme \ref{morphisms-lemma-composition-proper},
\item H-projectif, voir
Morphismes, Lemme \ref{morphisms-lemma-H-projective-composition},
\item fini ou entier, voir
Morphismes, Lemme \ref{morphisms-lemma-composition-finite},
\item fini localement libre, voir
Morphismes, Lemme \ref{morphisms-lemma-composition-finite-locally-free},
\item universellement submersif, voir
Morphismes, Lemme \ref{morphisms-lemma-composition-universally-submersive},
\item homéomorphisme universel, voir
Morphismes, Lemme \ref{morphisms-lemma-composition-universal-homeomorphism}.
\end{enumerate}
Compléter au besoin.
\end{remark}

\begin{remark}
\label{remark-list-properties-fpqc-local-base}
Parmi les propriétés mentionnées qui sont stables par changement de base
(énumérées dans la Remarque \ref{remark-list-properties-stable-base-change}),
les suivantes sont aussi {\it locales sur la base pour la topologie fpqc}
(et a fortiori locales sur la base pour la topologie fppf) :
\begin{enumerate}
\item pour les immersions, c'est le cas pour
\begin{enumerate}
\item les immersions fermées, voir
Descente, Lemme \ref{descent-lemma-descending-property-closed-immersion},
\item les immersions ouvertes, voir
Descente, Lemme \ref{descent-lemma-descending-property-open-immersion}, et
\item les immersions quasi-compactes, voir
Descente,
Lemme \ref{descent-lemma-descending-property-quasi-compact-immersion},
\end{enumerate}
\item quasi-compact, voir
Descente, Lemme \ref{descent-lemma-descending-property-quasi-compact},
\item universellement fermé, voir
Descente, Lemme
\ref{descent-lemma-descending-property-universally-closed},
\item (quasi-)séparé, voir
Descente, Lemmes
\ref{descent-lemma-descending-property-quasi-separated}, et
\ref{descent-lemma-descending-property-separated},
\item monomorphisme, voir
Descente, Lemme \ref{descent-lemma-descending-property-monomorphism},
\item surjectif, voir
Descente, Lemme \ref{descent-lemma-descending-property-surjective},
\item radiciel, voir
Descente, Lemme \ref{descent-lemma-descending-property-universally-injective},
\item affine, voir
Descente, Lemme \ref{descent-lemma-descending-property-affine},
\item quasi-affine, voir
Descente, Lemme \ref{descent-lemma-descending-property-quasi-affine},
\item (localement) de type fini, voir
Descente,
Lemmes \ref{descent-lemma-descending-property-locally-finite-type}, et
\ref{descent-lemma-descending-property-finite-type},
\item (localement) quasi-fini, voir
Descente, Lemme \ref{descent-lemma-descending-property-quasi-finite},
\item (localement) de présentation finie, voir
Descente, Lemmes
\ref{descent-lemma-descending-property-locally-finite-presentation}, et
\ref{descent-lemma-descending-property-finite-presentation},
\item localement de type fini de dimension relative $d$, voir
Descente,
Lemme \ref{descent-lemma-descending-property-relative-dimension-d},
\item universellement ouvert, voir
Descente, Lemme \ref{descent-lemma-descending-property-universally-open},
\item plat, voir
Descente, Lemme \ref{descent-lemma-descending-property-flat},
\item syntomique, voir
Descente, Lemme \ref{descent-lemma-descending-property-syntomic},
\item lisse, voir
Descente, Lemme \ref{descent-lemma-descending-property-smooth},
\item non ramifié (resp.\ G-non ramifié), voir
Descente, Lemme \ref{descent-lemma-descending-property-unramified},
\item \'etale, voir
Descente, Lemme \ref{descent-lemma-descending-property-etale},
\item propre, voir
Descente, Lemme \ref{descent-lemma-descending-property-proper},
\item fini ou entier, voir
Descente, Lemme \ref{descent-lemma-descending-property-finite},
\item fini localement libre, voir
Descente, Lemme \ref{descent-lemma-descending-property-finite-locally-free},
\item universellement submersif, voir
Descente, Lemme \ref{descent-lemma-descending-property-universally-submersive},
\item homéomorphisme universel, voir
Descente, Lemme \ref{descent-lemma-descending-property-universal-homeomorphism}.
\end{enumerate}
Notons que la propriété d'être une ``immersion'' n'est pas nécessairement locale
sur la base pour la topologie fpqc, mais que dans
Descente, Lemme \ref{descent-lemma-descending-fppf-property-immersion},
nous avons démontré qu'elle est locale sur la base pour la topologie fppf.
\end{remark}








\section{Propriétés des morphismes représentables de préfaisceaux}
\label{section-representable-properties}

\noindent
Voici la définition qui permet de procéder.

\begin{definition}
\label{definition-relative-representable-property}
On se donne $S$ et $a : F \to G$ représentable comme ci-dessus.
Soit $\mathcal{P}$ une propriété de morphismes de schémas qui
\begin{enumerate}
\item est stable par tout changement de base,
voir Schémas, Définition \ref{schemes-definition-preserved-by-base-change},
et
\item est locale sur la base pour la topologie fppf, voir
Descente, Définition \ref{descent-definition-property-morphisms-local}.
\end{enumerate}
Dans ce cas, nous disons que $a$ possède {\it la propriété $\mathcal{P}$} si, pour tout
$U \in \Ob((\Sch/S)_{fppf})$ et
tout $\xi \in G(U)$, le morphisme de schémas qui en résulte
$V_\xi \to U$ possède la propriété $\mathcal{P}$.
\end{definition}

\noindent
Il importe de noter que nous n'utiliserons cette définition que pour
des propriétés de morphismes stables par changement de base et
locales sur la base pour la topologie fppf. Ce choix ne vient
pas de ce que la définition n'aurait autrement aucun sens ; il vient plutôt
de ce que nous pouvons vouloir donner une autre définition
mieux adaptée à la propriété envisagée.

\begin{remark}
\label{remark-warning}
Considérons la propriété $\mathcal{P}=$``surjectif''.
Dans ce cas, dire
``soit $F \to G$ un morphisme surjectif'' pourrait être ambigu.
En effet, nous pourrions entendre la notion définie
dans la Définition \ref{definition-relative-representable-property}
ci-dessus, ou bien un morphisme surjectif de préfaisceaux, voir
Sites, Définition \ref{sites-definition-presheaves-injective-surjective},
ou encore, si $F$ et $G$ sont tous deux des faisceaux,
un morphisme surjectif de faisceaux, voir
Sites, Définition \ref{sites-definition-sheaves-injective-surjective}.
Sauf mention expresse du contraire, lorsque nous parlerons de morphismes d'espaces algébriques,
nous entendrons toujours la première notion. Voir le
Lemme \ref{lemma-surjective-flat-locally-finite-presentation}
pour un cas où cette propriété entraîne la surjectivité du morphisme de faisceaux sous-jacent.
\end{remark}

\noindent
Voici une vérification de cohérence.

\begin{lemma}
\label{lemma-morphism-schemes-gives-representable-transformation-property}
Soient $S$, $X$, $Y$ des objets de $\Sch_{fppf}$.
Soit $f : X \to Y$ un morphisme de schémas.
Soit $\mathcal{P}$ comme dans la
Définition \ref{definition-relative-representable-property}.
Alors $h_X \longrightarrow h_Y$ possède la propriété $\mathcal{P}$ si
et seulement si $f$ possède la propriété $\mathcal{P}$.
\end{lemma}

\begin{proof}
Notons que le lemme a un sens en vertu du
Lemme \ref{lemma-morphism-schemes-gives-representable-transformation}.
Démonstration omise.
\end{proof}

\begin{lemma}
\label{lemma-composition-representable-transformations-property}
Soit $S$ un schéma appartenant à $\Sch_{fppf}$.
Soient $F, G, H : (\Sch/S)_{fppf}^{opp} \to \textit{Ens}$.
Soit $\mathcal{P}$ une propriété comme dans la
Définition \ref{definition-relative-representable-property}
qui est stable par composition.
Soient $a : F \to G$, $b : G \to H$ des morphismes représentables de foncteurs.
Si $a$ et $b$ possèdent la propriété $\mathcal{P}$, il en va de même de
$b \circ a : F \longrightarrow H$.
\end{lemma}

\begin{proof}
Notons que le lemme a un sens en vertu du
Lemme \ref{lemma-composition-representable-transformations}.
Démonstration omise.
\end{proof}

\begin{lemma}
\label{lemma-base-change-representable-transformations-property}
Soit $S$ un schéma appartenant à $\Sch_{fppf}$.
Soient $F, G, H : (\Sch/S)_{fppf}^{opp} \to \textit{Ens}$.
Soit $\mathcal{P}$ une propriété comme dans la
Définition \ref{definition-relative-representable-property}.
Soit $a : F \to G$ un morphisme représentable de foncteurs.
Soit $b : H \to G$ un morphisme quelconque de foncteurs.
Considérons le diagramme de produit fibré
$$
\xymatrix{
H \times_{b, G, a} F \ar[r]_-{b'} \ar[d]_{a'} & F \ar[d]^a \\
H \ar[r]^b & G
}
$$
Si $a$ possède la propriété $\mathcal{P}$, alors son changement de base $a'$
possède également la propriété $\mathcal{P}$.
\end{lemma}

\begin{proof}
Notons que le lemme a un sens en vertu du
Lemme \ref{lemma-base-change-representable-transformations}.
Démonstration omise.
\end{proof}

\begin{lemma}
\label{lemma-descent-representable-transformations-property}
Soit $S$ un schéma appartenant à $\Sch_{fppf}$.
Soient $F, G, H : (\Sch/S)_{fppf}^{opp} \to \textit{Ens}$.
Soit $\mathcal{P}$ une propriété comme dans la
Définition \ref{definition-relative-representable-property}.
Soit $a : F \to G$ un morphisme représentable de foncteurs.
Soit $b : H \to G$ un morphisme quelconque de foncteurs.
Considérons le diagramme de produit fibré
$$
\xymatrix{
H \times_{b, G, a} F \ar[r]_-{b'} \ar[d]_{a'} & F \ar[d]^a \\
H \ar[r]^b & G
}
$$
Supposons que $b$ induise un morphisme surjectif de faisceaux pour la topologie fppf $H^\# \to G^\#$.
Dans ce cas, si $a'$ possède la propriété $\mathcal{P}$, alors $a$
possède également la propriété $\mathcal{P}$.
\end{lemma}

\begin{proof}
Remarquons d'abord que, d'après le
Lemme \ref{lemma-base-change-representable-transformations},
le morphisme $a'$ est représentable.
Soient $U \in \Ob((\Sch/S)_{fppf})$ et
$\xi \in G(U)$. Par hypothèse, il existe un recouvrement fppf
$\{U_i \to U\}_{i \in I}$ et des éléments $\xi_i \in H(U_i)$
qui sont envoyés sur $\xi|_{U_i}$ par $b$. Il résulte de la théorie générale des catégories que, pour
chaque $i$, on a un diagramme de produit fibré
$$
\xymatrix{
U_i \times_{\xi_i, H, a'} (H \times_{b, G, a} F) \ar[r] \ar[d] &
U \times_{\xi, G, a} F \ar[d] \\
U_i \ar[r] & U
}
$$
Par hypothèse, la flèche verticale de gauche est un morphisme de schémas qui
possède la propriété $\mathcal{P}$. Comme $\mathcal{P}$ est locale sur la base pour la
topologie fppf, il en résulte que la flèche verticale de droite possède aussi la propriété
$\mathcal{P}$, comme souhaité.
\end{proof}

\begin{lemma}
\label{lemma-product-representable-transformations-property}
Soit $S$ un schéma appartenant à $\Sch_{fppf}$.
Soient $F_i, G_i : (\Sch/S)_{fppf}^{opp} \to \textit{Ens}$,
$i = 1, 2$.
Soient $a_i : F_i \to G_i$, $i = 1, 2$, des morphismes représentables
de foncteurs.
Soit $\mathcal{P}$ une propriété comme dans la
Définition \ref{definition-relative-representable-property}
qui est stable par composition.
Si $a_1$ et $a_2$ possèdent la propriété $\mathcal{P}$, il en va de même de
$a_1 \times a_2 : F_1 \times F_2 \longrightarrow G_1 \times G_2$.
\end{lemma}

\begin{proof}
Notons que le lemme a un sens en vertu du
Lemme \ref{lemma-product-representable-transformations}.
Démonstration omise.
\end{proof}

\begin{lemma}
\label{lemma-representable-transformations-property-implication}
Soit $S$ un schéma appartenant à $\Sch_{fppf}$.
Soient $F, G : (\Sch/S)_{fppf}^{opp} \to \textit{Ens}$.
Soit $a : F \to G$ un morphisme représentable de foncteurs.
Soient $\mathcal{P}$, $\mathcal{P}'$ des propriétés comme dans la
Définition \ref{definition-relative-representable-property}.
Supposons que, pour tout morphisme de schémas $f : X \to Y$,
on ait $\mathcal{P}(f) \Rightarrow \mathcal{P}'(f)$.
Si $a$ possède la propriété $\mathcal{P}$, alors
$a$ possède la propriété $\mathcal{P}'$.
\end{lemma}

\begin{proof}
Formel.
\end{proof}

\begin{lemma}
\label{lemma-surjective-flat-locally-finite-presentation}
Soit $S$ un schéma.
Soient $F, G : (\Sch/S)_{fppf}^{opp} \to \textit{Ens}$ des faisceaux.
Soit $a : F \to G$ représentable, plat,
localement de présentation finie et surjectif.
Alors $a : F \to G$ est surjectif en tant que morphisme de faisceaux.
\end{lemma}

\begin{proof}
Soient $T$ un schéma au-dessus de $S$ et $g : T \to G$ un point à valeurs dans $T$ du foncteur
$G$. Par hypothèse, $T' = F \times_G T$ est (représenté par) un schéma et
le morphisme $T' \to T$ est plat, localement de présentation finie et
surjectif. Ainsi, $\{T' \to T\}$ est un recouvrement fppf tel
que $g|_{T'} \in G(T')$ provient d'un élément de $F(T')$, à savoir
le morphisme $T' \to F$. Cela prouve que le morphisme est surjectif en tant que
morphisme de faisceaux, voir
Sites, Définition \ref{sites-definition-sheaves-injective-surjective}.
\end{proof}

\noindent
Voici une caractérisation des foncteurs dont la
diagonale est représentable.

\begin{lemma}
\label{lemma-representable-diagonal}
Soit $S$ un schéma appartenant à $\Sch_{fppf}$.
Soit $F$ un préfaisceau d'ensembles sur $(\Sch/S)_{fppf}$.
Les assertions suivantes sont équivalentes :
\begin{enumerate}
\item la diagonale $F \to F \times F$ est représentable,
\item pour $U \in \Ob((\Sch/S)_{fppf})$ et tout $a \in F(U)$,
le morphisme $a : h_U \to F$ est représentable,
\item pour toute paire $U, V \in \Ob((\Sch/S)_{fppf})$
et tous $a \in F(U)$ et $b \in F(V)$, le produit fibré
$h_U \times_{a, F, b} h_V$ est représentable.
\end{enumerate}
\end{lemma}

\begin{proof}
C'est purement formel, voir
Catégories, Lemme \ref{categories-lemma-representable-diagonal}.
Cela ne dépend que du fait que la catégorie $(\Sch/S)_{fppf}$
possède les produits de deux objets et les produits fibrés, voir
Topologies, Lemme \ref{topologies-lemma-fibre-products-fppf}.
\end{proof}

\noindent
Dans la situation du lemme, pour tout morphisme
$\xi : h_U \to F$ comme dans son énoncé, il est légitime
de dire que $\xi$ possède la propriété $\mathcal{P}$, pour toute propriété
comme dans la Définition \ref{definition-relative-representable-property}.
C'est notamment le cas pour $\mathcal{P} = $ ``surjectif''
et $\mathcal{P} = $ ``\'etale'', voir la
Remarque \ref{remark-list-properties-fpqc-local-base}
ci-dessus. Nous utiliserons cette remarque dans la définition
des espaces algébriques ci-dessous.

\begin{lemma}
\label{lemma-transformation-diagonal-properties}
Soit $S$ un schéma appartenant à $\Sch_{fppf}$.
Soit $F$ un préfaisceau d'ensembles sur $(\Sch/S)_{fppf}$.
Soit $\mathcal{P}$ une propriété comme dans la
Définition \ref{definition-relative-representable-property}.
Supposons que, pour tous $U, V \in \Ob((\Sch/S)_{fppf})$, $a \in F(U)$ et
$b \in F(V)$, on ait
\begin{enumerate}
\item $h_U \times_{a, F, b} h_V$ est représentable, disons par le schéma $W$, et
\item le morphisme $W \to U \times_S V$ correspondant au morphisme
$h_U \times_{a, F, b} h_V \to h_U \times h_V$ possède la propriété $\mathcal{P}$,
\end{enumerate}
alors $\Delta : F \to F \times F$ est représentable et possède
la propriété $\mathcal{P}$.
\end{lemma}

\begin{proof}
Observons que $\Delta$ est représentable d'après le
Lemme \ref{lemma-representable-diagonal}.
Nous pouvons reformuler la condition (2) en disant que
le morphisme $h_U \times_{a, F, b} h_V \to h_{U \times_S V}$
possède la propriété $\mathcal{P}$, voir le Lemme
\ref{lemma-morphism-schemes-gives-representable-transformation-property}.
Considérons $T \in \Ob((\Sch/S)_{fppf})$ et $(a, b) \in (F \times F)(T)$.
Observons que nous avons le diagramme commutatif
$$
\xymatrix{
F \times_{\Delta, F \times F, (a, b)} h_T \ar[d] \ar[r] &
h_T \ar[d]^{\Delta_{T/S}} \\
h_T \times_{a, F, b} h_T \ar[r] \ar[d] &
h_{T \times_S T} \ar[d]^{(a, b)} \\
F \ar[r]^\Delta & F \times F
}
$$
dont les deux carrés sont cartésiens. Nous voyons ainsi que
le morphisme $F \times_{F \times F} h_T \to h_T$ est obtenu par changement
de base d'un morphisme qui possède la propriété $\mathcal{P}$ au moyen de
$\Delta_{T/S}$. Comme $\mathcal{P}$ est stable par changement de
base, cela achève la démonstration.
\end{proof}


















\section{Espaces algébriques}
\label{section-algebraic-spaces}

\noindent
Voici la définition.

\begin{definition}
\label{definition-algebraic-space}
Soit $S$ un schéma appartenant à $\Sch_{fppf}$.
Un {\it espace algébrique au-dessus de $S$} est un préfaisceau
$$
F : (\Sch/S)^{opp}_{fppf} \longrightarrow \textit{Ens}
$$
qui possède les propriétés suivantes :
\begin{enumerate}
\item Le préfaisceau $F$ est un faisceau.
\item Le morphisme diagonal $F  \to F \times F$ est représentable.
\item Il existe un schéma $U \in \Ob((\Sch/S)_{fppf})$
et un morphisme $h_U \to F$ qui est surjectif et \'etale\footnote{Voir le
Lemme \ref{lemma-representable-diagonal}, la
Définition \ref{definition-relative-representable-property} et la
Remarque \ref{remark-warning}.}.
\end{enumerate}
\end{definition}

\noindent
Notre définition diffère en deux points de la définition ``usuelle'', par exemple de la
définition donnée dans l'ouvrage de Knutson \cite{Kn}.

\medskip\noindent
La première est que nous exigeons que $F$ soit un faisceau pour la topologie fppf.
L'une des raisons de ce choix est que de nombreux exemples naturels
d'espaces algébriques vérifient la condition de faisceau pour les recouvrements fppf
(et même pour les recouvrements fpqc). En outre, l'utilité des espaces
algébriques tient notamment aux résultats de Michael Artin à leur sujet.
Sa méthode comporte une condition qui garantit que le résultat est
localement de présentation finie au-dessus de $S$.
Tout bien considéré, il nous semble que la topologie fppf
est la topologie naturelle à employer. Finalement, on obtient
la même catégorie d'espaces algébriques. Voir
Amorçage, section \ref{bootstrap-section-spaces-etale}.

\medskip\noindent
La seconde est que nous exigeons seulement que le morphisme diagonal de $F$ soit
représentable, tandis que \cite{Kn} exige qu'il soit également
quasi-compact. Si $F = h_U$ pour un schéma $U$ au-dessus de $S$,
cela revient à demander que $U$ soit quasi-séparé.
Nous cherchons à démontrer un certain
nombre des résultats qui suivent en supposant seulement que la diagonale
de $F$ soit représentable, et à ajouter simplement une hypothèse supplémentaire chaque fois que
cela est nécessaire. Cette convention a en tout cas pour conséquence agréable que
le lemme suivant est vrai.

\begin{lemma}
\label{lemma-scheme-is-space}
Tout schéma est un espace algébrique. Plus précisément,
pour tout schéma $T \in \Ob((\Sch/S)_{fppf})$,
le foncteur représentable $h_T$ est un espace algébrique.
\end{lemma}

\begin{proof}
Le foncteur $h_T$ est un faisceau d'après nos remarques de la section \ref{section-general}.
La diagonale $h_T \to h_T \times h_T = h_{T \times T}$ est
représentable parce que $(\Sch/S)_{fppf}$ possède les produits fibrés.
Le morphisme identité $h_T \to h_T$ est surjectif et \'etale.
\end{proof}

\begin{definition}
\label{definition-morphism-algebraic-spaces}
Soient $F$, $F'$ des espaces algébriques au-dessus de $S$.
Un {\it morphisme $f : F \to F'$ d'espaces algébriques au-dessus de $S$}
est un morphisme de foncteurs de $F$ vers $F'$.
\end{definition}

\noindent
La catégorie des espaces algébriques au-dessus de $S$ contient la catégorie
$(\Sch/S)_{fppf}$ comme sous-catégorie pleine par l'intermédiaire du
plongement de Yoneda $T/S \mapsto h_T$. Désormais, nous ne distinguerons plus
un schéma $T/S$ de l'espace algébrique qu'il représente.
Ainsi, lorsque nous disons ``Soit $f : T \to F$ un morphisme du schéma
$T$ vers l'espace algébrique $F$'', nous voulons dire que
$T \in \Ob((\Sch/S)_{fppf})$, que $F$ est un
espace algébrique au-dessus de $S$ et que $f : h_T \to F$ est un morphisme
d'espaces algébriques au-dessus de $S$.






\section{Produits fibrés d'espaces algébriques}
\label{section-fibre-products}

\noindent
La catégorie des espaces algébriques au-dessus de $S$ possède à la fois des produits et des
produits fibrés.

\begin{lemma}
\label{lemma-product-spaces}
Soit $S$ un schéma appartenant à $\Sch_{fppf}$.
Soient $F, G$ des espaces algébriques au-dessus de $S$.
Alors $F \times G$ est un espace algébrique et c'est un produit
dans la catégorie des espaces algébriques au-dessus de $S$.
\end{lemma}

\begin{proof}
Il est immédiat que $H = F \times G$ est un faisceau.
La diagonale de $H$ est simplement le produit des
diagonales de $F$ et de $G$. Elle est donc représentable d'après le
Lemme \ref{lemma-product-representable-transformations}.
Enfin, si $U \to F$ et $V \to G$ sont des morphismes étales
surjectifs, où $U, V \in \Ob((\Sch/S)_{fppf})$,
alors $U \times V \to F \times G$ est étale surjectif
d'après le Lemme \ref{lemma-product-representable-transformations-property}.
\end{proof}

\begin{lemma}
\label{lemma-fibre-product-spaces-over-sheaf-with-representable-diagonal}
Soit $S$ un schéma appartenant à $\Sch_{fppf}$.
Soit $H$ un faisceau sur $(\Sch/S)_{fppf}$ dont la diagonale
est représentable. Soient $F, G$ des espaces algébriques au-dessus de $S$.
Soient $F \to H$, $G \to H$ des morphismes de faisceaux.
Alors $F \times_H G$ est un espace algébrique.
\end{lemma}

\begin{proof}
Vérifions les trois conditions de la
Définition \ref{definition-algebraic-space}.
Un produit fibré de faisceaux est un faisceau; ainsi, $F \times_H G$ est un faisceau.
La diagonale de $F \times_H G$ est la flèche verticale de gauche dans le diagramme
$$
\xymatrix{
F \times_H G \ar[r] \ar[d]_\Delta &
F \times G \ar[d]^{\Delta_F \times \Delta_G} \\
(F \times F) \times_{(H \times H)} (G \times G) \ar[r] &
(F \times F) \times (G \times G)
}
$$
qui est cartésien. Ainsi, $\Delta$ est représentable, puisqu'elle s'obtient par changement de base
à partir du morphisme de droite, qui est représentable; voir les
Lemmes \ref{lemma-product-representable-transformations} et
\ref{lemma-base-change-representable-transformations}.
Enfin, soient $U, V \in \Ob((\Sch/S)_{fppf})$
et soient $a : U \to F$, $b : V \to G$ des morphismes étales surjectifs.
Comme $\Delta_H$ est représentable, on voit que $U \times_H V$ est un schéma.
Le morphisme
$$
U \times_H V \longrightarrow F \times_H G
$$
est étale surjectif, car il est le composé des morphismes
$U \times_H V \to U \times_H G$ et $U \times_H G \to F \times_H G$,
obtenus par changement de base à partir des morphismes étales surjectifs $U \to F$ et $V \to G$; voir les
Lemmes \ref{lemma-composition-representable-transformations} et
\ref{lemma-base-change-representable-transformations}.
Cela montre que la dernière condition de la
Définition \ref{definition-algebraic-space}
est satisfaite, et l'on conclut que $F \times_H G$ est un espace algébrique.
\end{proof}

\begin{lemma}
\label{lemma-fibre-product-spaces}
Soit $S$ un schéma appartenant à $\Sch_{fppf}$.
Soient $F \to H$, $G \to H$ des morphismes d'espaces algébriques au-dessus de $S$.
Alors $F \times_H G$ est un espace algébrique et c'est un produit fibré
dans la catégorie des espaces algébriques au-dessus de $S$.
\end{lemma}

\begin{proof}
Il résulte du Lemme plus général
\ref{lemma-fibre-product-spaces-over-sheaf-with-representable-diagonal}
que $F \times_H G$ est un espace algébrique.
Il est clair que $F \times_H G$
est un produit fibré dans la catégorie des espaces algébriques au-dessus de $S$,
puisque celle-ci est une sous-catégorie pleine de la catégorie
des (pré)faisceaux d'ensembles sur $(\Sch/S)_{fppf}$.
\end{proof}






\section{Recollement d'espaces algébriques}
\label{section-glueing-algebraic-spaces}

\noindent
Dans cette section, nous commençons véritablement à abuser des notations et à ne plus
distinguer les schémas des espaces qu'ils représentent.

\begin{lemma}
\label{lemma-coproduct-sheaves-open-and-closed}
Soit $S \in \Ob(\Sch_{fppf})$. Soient $F$ et $G$ des faisceaux sur
$(\Sch/S)_{fppf}^{opp}$ et notons $F \amalg G$ leur somme
dans la catégorie des faisceaux. Le morphisme $F \to F \amalg G$ est représentable par
des immersions à la fois ouvertes et fermées.
\end{lemma}

\begin{proof}
Soient $U$ un schéma et $\xi \in (F \amalg G)(U)$. Rappelons que la
somme dans la catégorie des faisceaux est le faisceau associé au
préfaisceau somme (Sites, Lemme \ref{sites-lemma-colimit-sheaves}).
Il existe donc un recouvrement fppf $\{g_i : U_i \to U\}_{i \in I}$
et une décomposition en somme disjointe $I = I' \amalg I''$ telle que
$U_i \to U \to F \amalg G$ se factorise par $F$, resp.\ par $G$,
si et seulement si $i \in I'$, resp.\ $i \in I''$. Puisque l'intersection de $F$ et de
$G$ dans $F \amalg G$ est vide, on en déduit que
$U_i \times_U U_j$ est vide si $i \in I'$ et $j \in I''$.
Par conséquent, $U' = \bigcup_{i \in I'} g_i(U_i)$ et
$U'' = \bigcup_{i \in I''} g_i(U_i)$ sont des sous-schémas ouverts
disjoints de $U$ (Morphismes, Lemme \ref{morphisms-lemma-fppf-open})
tels que $U = U' \amalg U''$.
Nous omettons la vérification de l'égalité $U' = U \times_{F \amalg G} F$.
\end{proof}

\begin{lemma}
\label{lemma-representable-sheaf-coproduct-sheaves}
Soit $S \in \Ob(\Sch_{fppf})$.
Soit $U \in \Ob((\Sch/S)_{fppf})$.
Étant donnés un ensemble $I$ et des faisceaux $F_i$ sur $\Ob((\Sch/S)_{fppf})$,
si $U \cong \coprod_{i\in I} F_i$
comme faisceaux, chaque $F_i$ est représentable par un sous-schéma
$U_i$ à la fois ouvert et fermé, et $U \cong \coprod U_i$ comme schémas.
\end{lemma}

\begin{proof}
D'après le Lemme \ref{lemma-coproduct-sheaves-open-and-closed},
le morphisme $F_i \to U$ est représentable par des immersions à la fois ouvertes et fermées.
Ainsi, $F_i$ est représentable par un sous-schéma $U_i$ à la fois ouvert et fermé de $U$.
On a $U = \coprod U_i$, puisque $U \cong \coprod F_i$
comme faisceaux et que l'égalité peut être vérifiée sur les points.
\end{proof}

\begin{lemma}
\label{lemma-algebraic-space-coproduct-sheaves}
Soit $S \in \Ob(\Sch_{fppf})$.
Soit $F$ un espace algébrique au-dessus de $S$.
Étant donnés un ensemble $I$ et des faisceaux $F_i$ sur
$\Ob((\Sch/S)_{fppf})$,
si $F \cong \coprod_{i\in I} F_i$ comme faisceaux,
chaque $F_i$ est un espace algébrique au-dessus de $S$.
\end{lemma}

\begin{proof}
La représentabilité de $F \to F \times F$ implique que chacun des morphismes diagonaux
$F_i \to F_i \times F_i$ est représentable (cela résulte immédiatement des définitions
et de l'égalité $F \times_{(F \times F)} (F_i \times F_i) = F_i$).
Choisissons un schéma $U$ dans $(\Sch/S)_{fppf}$ et un morphisme
\'etale surjectif $U \to F$ (dont l'existence résulte de l'hypothèse).
Le morphisme $U \times_F F_i \to F_i$ obtenu par changement de base est \'etale surjectif
d'après le Lemme \ref{lemma-base-change-representable-transformations-property}.
D'autre part, $U \times_F F_i$ est un schéma d'après le
Lemme \ref{lemma-coproduct-sheaves-open-and-closed}.
Nous avons donc vérifié toutes les conditions de la
Définition \ref{definition-algebraic-space},
et $F_i$ est un espace algébrique.
\end{proof}

\noindent
La condition portant sur la taille de $I$ et des $F_i$ dans le
lemme suivant peut être ignorée par ceux que ne préoccupent pas les
questions de théorie des ensembles.

\begin{lemma}
\label{lemma-coproduct-algebraic-spaces}
Soit $S \in \Ob(\Sch_{fppf})$.
Supposons donnés un ensemble $I$ et des espaces algébriques $F_i$, $i \in I$.
Alors $F = \coprod_{i \in I} F_i$ est un espace algébrique,
pourvu que $I$ et les $F_i$ ne soient pas trop «~grands~» : par exemple, si nous
pouvons choisir des morphismes \'etales surjectifs $U_i \to F_i$ tels que
$\coprod_{i \in I} U_i$ soit isomorphe à un objet de
$(\Sch/S)_{fppf}$, alors $F$ est un espace algébrique.
\end{lemma}

\begin{proof}
Par construction, $F$ est un faisceau. Nous omettons de vérifier que le
morphisme diagonal de $F$ est représentable. Enfin, si $U$ est un
objet de $(\Sch/S)_{fppf}$ isomorphe à $\coprod_{i \in I} U_i$,
alors il est immédiat de vérifier que le morphisme qui en résulte
$U \to \coprod F_i$ est \'etale surjectif.
\end{proof}

\noindent
Voici l'analogue de Schémas, Lemme \ref{schemes-lemma-glue-functors}.

\begin{lemma}
\label{lemma-glueing-algebraic-spaces}
Soit $S \in \Ob(\Sch_{fppf})$.
Soit $F$ un préfaisceau d'ensembles sur $(\Sch/S)_{fppf}$.
Supposons que
\begin{enumerate}
\item $F$ est un faisceau,
\item il existe un ensemble d'indices $I$
et des sous-foncteurs $F_i \subset F$ tels que
\begin{enumerate}
\item chaque $F_i$ est un espace algébrique,
\item chaque $F_i \to F$ est représentable,
\item chaque $F_i \to F$ est une immersion ouverte (voir la
Définition \ref{definition-relative-representable-property}),
\item le morphisme $\coprod F_i \to F$ est surjectif comme morphisme de faisceaux, et
\item $\coprod F_i$ est un espace algébrique (condition ensembliste, voir le
Lemme \ref{lemma-coproduct-algebraic-spaces}).
\end{enumerate}
\end{enumerate}
Alors $F$ est un espace algébrique.
\end{lemma}

\begin{proof}
Soit $T$ un objet de $(\Sch/S)_{fppf}$. Soit $T \to F$ un morphisme.
D'après les hypothèses (2)(b) et (2)(c), le produit fibré $F_i \times_F T$
est représentable par un sous-schéma ouvert $V_i \subset T$. Il s'ensuit que
$(\coprod F_i) \times_F T$ est représenté par le schéma $\coprod V_i$ au-dessus de $T$.
D'après l'hypothèse (2)(d), il existe un recouvrement fppf $\{T_j \to T\}_{j \in J}$
tel que $T_j \to T \to F$ se factorise par $F_i$, où $i = i(j)$.
Ainsi, $T_j \to T$ se factorise par le sous-schéma ouvert $V_{i(j)} \subset T$.
Puisque la famille $\{T_j \to T\}$ est surjective, il s'ensuit que
$T = \bigcup V_i$ est un recouvrement ouvert. En particulier, le morphisme
de foncteurs $\coprod F_i \to F$ est représentable
et surjectif au sens de la
Définition \ref{definition-relative-representable-property}
(voir la Remarque \ref{remark-warning} pour une discussion).

\medskip\noindent
Ensuite, soit $T' \to F$ un second morphisme ayant pour source un objet de $(\Sch/S)_{fppf}$.
Écrivons comme ci-dessus $T' = \bigcup V'_i$, où $V'_i = T' \times_F F_i$.
Pour montrer que la diagonale $F \to F \times F$ est représentable,
il faut montrer que $G = T \times_F T'$ est représentable; voir le
Lemme \ref{lemma-representable-diagonal}.
Considérons les sous-foncteurs $G_i = G \times_F F_i$.
Notons que $G_i = V_i \times_{F_i} V'_i$; il est donc représentable,
car $F_i$ est un espace algébrique.
D'après ce qui précède, les $G_i$ forment un recouvrement de Zariski de $G$.
Ainsi, d'après Schémas, Lemme \ref{schemes-lemma-glue-functors},
on voit que $G$ est représentable.

\medskip\noindent
Choisissons un schéma $U \in \Ob((\Sch/S)_{fppf})$ et un morphisme
\'etale surjectif $U \to \coprod F_i$ (dont l'existence résulte de l'hypothèse).
Nous pouvons écrire $U = \coprod U_i$, où $U_i$ est l'image réciproque de $F_i$;
voir le Lemme \ref{lemma-representable-sheaf-coproduct-sheaves}.
Nous affirmons que $U \to F$ est \'etale surjectif. La surjectivité découle
de celle de $\coprod F_i \to F$ (voir le premier paragraphe de la démonstration),
en appliquant le
Lemme \ref{lemma-composition-representable-transformations-property}.
Considérons le produit fibré $U \times_F T$, où $T \to F$ est comme
ci-dessus. Il faut montrer que $U \times_F T \to T$ est \'etale.
Puisque $U \times_F T = \coprod U_i \times_F T$, il suffit de montrer que
chacun des morphismes $U_i \times_F T \to T$ est \'etale. Comme
$U_i \times_F T = U_i \times_{F_i} V_i$, cela résulte du
fait que $U_i \to F_i$ est \'etale et que $V_i \to T$ est une immersion ouverte
(voir aussi Morphismes, Lemmes \ref{morphisms-lemma-open-immersion-etale}
et \ref{morphisms-lemma-composition-etale}).
\end{proof}














\section{Présentations d'espaces algébriques}
\label{section-presentations}

\noindent
Étant donné un espace algébrique, on peut en trouver une «~présentation~».


\begin{lemma}
\label{lemma-space-presentation}
Soit $F$ un espace algébrique au-dessus de $S$. Soit $f : U \to F$ un
morphisme \'etale surjectif d'un schéma vers $F$. Posons $R = U \times_F U$.
Alors
\begin{enumerate}
\item $j : R \to U \times_S U$ définit sur
$U$ une relation d'équivalence au-dessus de $S$ (voir
Groupoïdes, Définition \ref{groupoids-definition-equivalence-relation}).
\item les morphismes $s, t : R \to U$ sont \'etales, et
\item le diagramme
$$
\xymatrix{
R \ar@<1ex>[r] \ar@<-1ex>[r] &
U \ar[r] &
F
}
$$
est un diagramme de conoyau dans $\Sh((\Sch/S)_{fppf})$.
\end{enumerate}
\end{lemma}

\begin{proof}
Soit $T/S$ un objet de $(\Sch/S)_{fppf}$.
Alors $R(T) = \{(a, b) \in U(T) \times U(T) \mid f \circ a = f \circ b\}$,
ce qui définit une relation d'équivalence sur $U(T)$.
Les morphismes $s, t : R \to U$ sont \'etales parce que le morphisme
$U \to F$ est \'etale.

\medskip\noindent
Pour démontrer (3), montrons d'abord que
$U \to F$ est un morphisme surjectif de faisceaux; voir
Sites, Définition \ref{sites-definition-sheaves-injective-surjective}.
Soit $\xi \in F(T)$, où $T$ est comme ci-dessus. Posons $V = T \times_{\xi, F, f}U$.
Par hypothèse, $V$ est un schéma et $V \to T$ est un morphisme \'etale surjectif.
Ainsi, $\{V \to T\}$ est un recouvrement pour la topologie fppf.
Puisque, par construction, $\xi|_V$ se factorise par $U$, on
en déduit que $U \to F$ est surjectif. La surjectivité implique que
$F$ est le conoyau du diagramme d'après
Sites, Lemme \ref{sites-lemma-coequalizer-surjection}.
\end{proof}

\noindent
Ce lemme conduit aux définitions suivantes.

\begin{definition}
\label{definition-etale-equivalence-relation}
Soit $S$ un schéma. Soit $U$ un schéma au-dessus de $S$.
Une {\it relation d'équivalence \'etale} sur $U$ au-dessus de $S$
est une relation d'équivalence $j : R \to U \times_S U$
telle que $s, t : R \to U$ soient des morphismes \'etales de schémas.
\end{definition}

\begin{definition}
\label{definition-presentation}
Soit $F$ un espace algébrique au-dessus de $S$.
Une {\it présentation} de $F$ est la donnée d'un schéma
$U$ au-dessus de $S$, d'une relation d'équivalence \'etale $R$ sur $U$ au-dessus de $S$, et
d'un morphisme \'etale surjectif $U \to F$ tel que $R = U \times_F U$.
\end{definition}

\noindent
De manière équivalente, on pourrait demander l'existence d'un isomorphisme
$$
U/R \cong F
$$
où le quotient $U/R$ est défini comme dans
Groupoïdes, section \ref{groupoids-section-quotient-sheaves}.
Pour construire des espaces algébriques, nous étudierons la question réciproque, à savoir
pour quelles relations d'équivalence le faisceau quotient $U/R$ est un espace algébrique.
On verra finalement qu'il en est toujours ainsi si $R$ est une relation
d'équivalence \'etale sur $U$ au-dessus de $S$; voir le Théorème \ref{theorem-presentation}.





























\section{Espaces algébriques et relations d'équivalence}
\label{section-spaces-from-equivalence-relations}

\noindent
Supposons donnés un schéma $U$ au-dessus de $S$
et une relation d'équivalence \'etale $R$ sur $U$ au-dessus de $S$.
Nous voudrions montrer que ces données définissent un espace algébrique.
Nous allons établir une suite de lemmes montrant que le faisceau quotient $U/R$
(voir Groupoïdes, Définition \ref{groupoids-definition-quotient-sheaf})
possède toutes les propriétés que lui impose la
Définition \ref{definition-algebraic-space}.

\begin{lemma}
\label{lemma-pullback-etale-equivalence-relation}
Soit $S$ un schéma. Soit $U$ un schéma au-dessus de $S$.
Soit $j = (s, t) : R \to U \times_S U$
une relation d'équivalence \'etale sur $U$ au-dessus de $S$.
Soit $U' \to U$ un morphisme \'etale.
Soit $R'$ la restriction de $R$ à $U'$; voir
Groupoïdes, Définition \ref{groupoids-definition-restrict-relation}.
Alors $j' : R' \to U' \times_S U'$ est également une relation
d'équivalence \'etale.
\end{lemma}

\begin{proof}
Il résulte de la description de $s', t'$ dans
Groupoïdes, Lemme \ref{groupoids-lemma-restrict-groupoid}
que $s' , t' : R' \to U'$ sont \'etales,
car ce sont des composés de changements de base de morphismes \'etales
(voir Morphismes, Lemme \ref{morphisms-lemma-base-change-etale}
et \ref{morphisms-lemma-composition-etale}).
\end{proof}

\noindent
Nous utiliserons souvent le lemme suivant pour trouver des sous-espaces ouverts des espaces
algébriques. Une légère amélioration de ce lemme (sous des hypothèses plus générales) est
Amorçage, Lemme \ref{bootstrap-lemma-better-finding-opens}.

\begin{lemma}
\label{lemma-finding-opens}
Soit $S$ un schéma.
Soit $U$ un schéma au-dessus de $S$.
Soit $j = (s, t) : R \to U \times_S U$ une pré-relation.
Soit $g : U' \to U$ un morphisme.
Supposons que
\begin{enumerate}
\item $j$ est une relation d'équivalence,
\item $s, t : R \to U$ sont surjectifs, plats et
localement de présentation finie,
\item $g$ est plat et localement de présentation finie.
\end{enumerate}
Soit $R' = R|_{U'}$ la restriction de $R$ à $U'$. Alors
$U'/R' \to U/R$ est représentable et est une immersion ouverte.
\end{lemma}

\begin{proof}
D'après Groupoïdes, Lemme \ref{groupoids-lemma-restrict-relation},
le morphisme $j' = (s', t') : R' \to U' \times_S U'$
définit une relation d'équivalence. Puisque $g$ est plat et localement de
présentation finie, $g$ est aussi universellement ouvert
(Morphismes, Lemme \ref{morphisms-lemma-fppf-open}).
Pour la même raison, $s, t$ sont également universellement ouverts.
Posons $W^1 = g(U') \subset U$ et $W = t(s^{-1}(W^1))$.
Alors $W^1$ et $W$ sont ouverts dans $U$. De plus, puisque $j$ est une
relation d'équivalence, on a $t(s^{-1}(W)) = W$ (voir
Groupoïdes, Lemme \ref{groupoids-lemma-constructing-invariant-opens},
par exemple).

\medskip\noindent
D'après
Groupoïdes,
Lemme \ref{groupoids-lemma-quotient-pre-equivalence-relation-restrict},
le morphisme de faisceaux $F' = U'/R' \to F = U/R$ est injectif.
Soit $a : T \to F$ un morphisme d'un schéma vers $U/R$.
Il faut montrer que $T \times_F F'$ est représentable
par un sous-schéma ouvert de $T$.

\medskip\noindent
Le morphisme $a$ est donné par les données suivantes :
un recouvrement fppf $\{\varphi_j : T_j \to T\}_{j \in J}$ de $T$ et des
morphismes $a_j : T_j \to U$ tels que les morphismes
$$
a_j \times a_{j'} :
T_j \times_T T_{j'}
\longrightarrow
U \times_S U
$$
se factorisent par $j : R \to U \times_S U$ au moyen de morphismes (uniques)
$r_{jj'} : T_j \times_T T_{j'} \to R$. Le système
$(a_j)$ correspond à $a$ en ce sens que les diagrammes
$$
\xymatrix{
T_j \ar[r]_{a_j} \ar[d] & U \ar[d] \\
T \ar[r]^a & F
}
$$
commutent.

\medskip\noindent
Considérons les ouverts $W_j = a_j^{-1}(W) \subset T_j$.
Puisque $t(s^{-1}(W)) = W$, on voit que
$$
W_j \times_T T_{j'} =
r_{jj'}^{-1}(t^{-1}(W)) = r_{jj'}^{-1}(s^{-1}(W)) =
T_j \times_T W_{j'}.
$$
D'après
Descente, Lemme \ref{descent-lemma-open-fpqc-covering},
cela signifie qu'il existe un ouvert
$W_T \subset T$ tel que $\varphi_j^{-1}(W_T) = W_j$ pour tout $j \in J$.
Nous affirmons que $W_T \to T$ représente $T \times_F F' \to T$.

\medskip\noindent
Montrons d'abord que $W_T \to T \to F$ définit un élément de
$F'(W_T)$. Puisque $\{W_j \to W_T\}_{j \in J}$ est un
recouvrement fppf de $W_T$, il suffit de montrer que
chaque $W_j \to U \to F$ définit un élément de $F'(W_j)$ (car $F'$ est un faisceau
pour la topologie fppf). Considérons le diagramme commutatif
$$
\xymatrix{
W'_j \ar[rr] \ar[dd] \ar[rd] & & U' \ar[d]^g \\
& s^{-1}(W^1) \ar[r]_s \ar[d]^t & W^1 \ar[d] \\
W_j \ar[r]^{a_j|_{W_j}} & W \ar[r] & F
}
$$
où $W'_j = W_j \times_W s^{-1}(W^1) \times_{W^1} U'$.
Puisque $t$ et $g$ sont surjectifs, plats et localement de présentation
finie, il en est de même de $W'_j \to W_j$. Par conséquent, la restriction de
l'élément $W_j \to U \to F$ à $W'_j$ est un élément de $F'$,
comme voulu.

\medskip\noindent
Supposons que $f : T' \to T$ soit un morphisme de schémas
tel que $a|_{T'} \in F'(T')$. Il faut montrer que
$f$ se factorise par l'ouvert $W_T$. Puisque
$\{T' \times_T T_j \to T'\}$ est un recouvrement fppf de $T'$,
il suffit de montrer que chaque $T' \times_T T_j \to T$
se factorise par $W_T$. Nous pouvons donc supposer que $f$ se factorise
sous la forme $\varphi_j \circ f_j : T' \to T_j \to T$ pour un certain $j$.
Dans ce cas, la condition $a|_{T'} \in F'(T')$ signifie qu'il existe
un recouvrement fppf $\{\psi_i : T'_i \to T'\}_{i \in I}$ et des
morphismes $b_i : T'_i \to U'$ tels que
$$
\xymatrix{
T'_i \ar[r]_{b_i} \ar[d]_{f_j \circ \psi_i} & U' \ar[r]_g & U \ar[d] \\
T_j \ar[r]^{a_j} & U \ar[r] & F
}
$$
soit commutatif. Cette commutativité signifie qu'il existe un
morphisme $r'_i : T'_i \to R$ tel que
$t \circ r'_i = a_j \circ f_j \circ \psi_i$ et
$s \circ r'_i = g \circ b_i$. Il en résulte que
$\Im(f_j \circ \psi_i) \subset W_j$, ce qui conclut.
\end{proof}

\noindent
Le lemme suivant n'est pas tout à fait trivial, bien qu'il semble
devoir l'être.

\begin{lemma}
\label{lemma-when-it-works-it-works}
Soit $S$ un schéma. Soit $U$ un schéma au-dessus de $S$.
Soit $j = (s, t) : R \to U \times_S U$
une relation d'équivalence \'etale sur $U$ au-dessus de $S$.
Si le quotient $U/R$ est un espace algébrique, alors
$U \to U/R$ est \'etale et surjectif. Par conséquent,
$(U, R, U \to U/R)$ est une présentation de l'espace
algébrique $U/R$.
\end{lemma}

\begin{proof}
Notons $c : U \to U/R$ le morphisme en question.
Soit $T$ un schéma et soit $a : T \to U/R$ un morphisme.
Il faut montrer que le morphisme (de schémas)
$\pi : T \times_{a, U/R, c} U \to T$ est \'etale et surjectif.
Le morphisme $a$ correspond à un recouvrement fppf
$\{\varphi_i : T_i \to T\}$ et à des morphismes $a_i : T_i \to U$ tels
que
$a_i \times a_{i'} : T_i \times_T T_{i'} \to U \times_S U$
se factorise par $R$, et tels que $c \circ a_i = a \circ \varphi_i$.
Ainsi,
$$
T_i \times_{\varphi_i, T} T \times_{a, U/R, c} U =
T_i \times_{c \circ a_i, U/R, c} U =
T_i \times_{a_i, U} U \times_{c, U/R, c} U = T_i \times_{a_i, U, t} R.
$$
Puisque $t$ est \'etale et surjectif, on en déduit que
le changement de base de $\pi$ à $T_i$ est surjectif et \'etale.
Comme la propriété d'être surjectif et \'etale est locale
sur la base pour la topologie fpqc (voir la
Remarque \ref{remark-list-properties-fpqc-local-base}),
cela conclut.
\end{proof}

\begin{lemma}
\label{lemma-presentation-quasi-compact}
Soit $S$ un schéma.
Soit $U$ un schéma au-dessus de $S$.
Soit $j = (s, t) : R \to U \times_S U$
une relation d'équivalence \'etale sur $U$ au-dessus de $S$.
Supposons $U$ affine. Alors le quotient $F = U/R$
est un espace algébrique, et $U \to F$ est \'etale et surjectif.
\end{lemma}

\begin{proof}
Puisque $j : R \to U \times_S U$ est un monomorphisme, $j$ est séparé
(voir Schémas, Lemme \ref{schemes-lemma-monomorphism-separated}).
Comme $U$ est affine, on voit que $U \times_S U$
(qui est muni d'un monomorphisme vers le schéma affine
$U \times U$) est séparé. Il s'ensuit que $R$ est séparé.
En particulier, les morphismes $s, t$ sont séparés et \'etales.

\medskip\noindent
Puisque le composé $R \to U \times_S U \to U$ est
localement de type fini, on en déduit que
$j$ est localement de type fini (voir
Morphismes, Lemme \ref{morphisms-lemma-permanence-finite-type}).
Comme $j$ est aussi un monomorphisme, ses fibres sont finies, et
l'on voit que $j$ est localement quasi-fini d'après
Morphismes, Lemme \ref{morphisms-lemma-finite-fibre}.
En définitive, $j$ est séparé et localement quasi-fini.

\medskip\noindent
La première étape consiste à montrer que le morphisme de passage au quotient
$c : U \to F$ est représentable.
Considérons un schéma $T$ et un morphisme $a : T \to F$.
Il faut montrer que le faisceau $G = T \times_{a, F, c} U$
est représentable.
Comme on l'a vu dans les démonstrations des Lemmes \ref{lemma-finding-opens} et
\ref{lemma-when-it-works-it-works}, il existe un recouvrement fppf
$\{\varphi_i : T_i \to T\}_{i \in I}$ et des morphismes $a_i : T_i \to U$
tels que $a_i \times a_{i'} : T_i \times_T T_{i'} \to U \times_S U$
se factorise par $R$, et tels que $c \circ a_i = a \circ \varphi_i$.
Comme dans la démonstration du Lemme \ref{lemma-when-it-works-it-works}, on voit que
\begin{eqnarray*}
T_i \times_{\varphi_i, T} G & = &
T_i \times_{\varphi_i, T} T \times_{a, U/R, c} U \\
& = & T_i \times_{c \circ a_i, U/R, c} U \\
& = & T_i \times_{a_i, U} U \times_{c, U/R, c} U \\
& = & T_i \times_{a_i, U, t} R
\end{eqnarray*}
Puisque $t$ est séparé et \'etale, donc en particulier
séparé et localement quasi-fini (d'après Morphismes, Lemmes
\ref{morphisms-lemma-unramified-quasi-finite} et
\ref{morphisms-lemma-flat-unramified-etale}),
on voit que la restriction
de $G$ à chaque $T_i$ est représentable par un morphisme de schémas
$X_i \to T_i$ séparé et localement quasi-fini. D'après
Descente, Lemme \ref{descent-lemma-descent-data-sheaves},
on obtient une donnée de descente $(X_i, \varphi_{ii'})$ relativement
au recouvrement fppf $\{T_i \to T\}$. Comme chaque
$X_i \to T_i$ est séparé et localement quasi-fini, il résulte de
Compléments sur les morphismes, Lemme
\ref{more-morphisms-lemma-separated-locally-quasi-finite-morphisms-fppf-descend}
que cette donnée de descente est effective.
Ainsi, d'après
Descente, Lemme \ref{descent-lemma-descent-data-sheaves} (2),
on en déduit que $G$ est représentable, comme voulu.

\medskip\noindent
La deuxième étape de la démonstration consiste à montrer que $U \to F$ est surjectif et
\'etale. Cela résulte de ce qui précède : dans la première étape, nous avons
vu que $G = T \times_{a, F, c} U$ est un schéma au-dessus de $T$ dont les changements de base
sont les schémas $X_i \to T_i$, qui sont surjectifs et \'etales. Ainsi, $G \to T$
est surjectif et \'etale (voir la
Remarque \ref{remark-list-properties-fpqc-local-base}).
On peut également reprendre la démonstration du
Lemme \ref{lemma-when-it-works-it-works} dans la présente
situation.

\medskip\noindent
La troisième et dernière étape consiste à montrer que le morphisme diagonal $F \to F \times F$
est représentable. Observons d'abord que le diagramme
$$
\xymatrix{
R \ar[r] \ar[d]_j & F \ar[d]^\Delta \\
U \times_S U \ar[r] & F \times F
}
$$
est un carré cartésien. D'après le
Lemme \ref{lemma-product-representable-transformations}, le morphisme
$U \times_S U \to F \times F$ est représentable (notons que
$h_U \times h_U = h_{U \times_S U}$). De plus, d'après le
Lemme \ref{lemma-product-representable-transformations-property},
le morphisme $U \times_S U \to F \times F$ est surjectif
et \'etale (notons aussi que \'etale et surjectif figurent dans les listes des
Remarques \ref{remark-list-properties-fpqc-local-base}
et \ref{remark-list-properties-stable-composition}).
Il résulte soit du
Lemme \ref{lemma-base-change-representable-transformations}
et du diagramme ci-dessus, soit de l'écriture de $R \to F$ comme $R \to U \to F$ et des
Lemmes
\ref{lemma-morphism-schemes-gives-representable-transformation} et
\ref{lemma-composition-representable-transformations}, que
$R \to F$ est également représentable. Soient $T$ un schéma et
$a : T \to F \times F$ un morphisme. Il faut montrer que
$G = T \times_{a, F \times F, \Delta} F$ est représentable.
D'après ce qui précède, le morphisme (de schémas)
$$
T' = (U \times_S U) \times_{F \times F, a} T \longrightarrow T
$$
est surjectif et \'etale. Ainsi, $\{T' \to T\}$ est un recouvrement
\'etale de $T$. Notons également que
$$
T' \times_T G = T' \times_{U \times_S U, j} R
$$
comme on le voit en examinant le cube suivant
$$
\xymatrix{
& R \ar[rr] \ar[dd] & & F \ar[dd] \\
T' \times_T G \ar[rr] \ar[dd] \ar[ru] & & G \ar[dd] \ar[ru] & \\
& U \times_S U \ar'[r][rr] & & F \times F \\
T' \ar[rr] \ar[ru] & & T \ar[ru]
}
$$
On voit ainsi que la restriction de $G$ à $T'$ est représentable
par un schéma $X$, et de plus que le morphisme $X \to T'$ est
un changement de base du morphisme $j$. Par conséquent, $X \to T'$ est
séparé et localement quasi-fini (voir le deuxième paragraphe de la démonstration).
D'après Descente, Lemme \ref{descent-lemma-descent-data-sheaves},
on obtient une donnée de descente $(X, \varphi)$ relativement
au recouvrement fppf $\{T' \to T\}$. Puisque
$X \to T'$ est séparé et localement quasi-fini, il résulte de
Compléments sur les morphismes, Lemme
\ref{more-morphisms-lemma-separated-locally-quasi-finite-morphisms-fppf-descend}
que cette donnée de descente est effective.
Ainsi, d'après
Descente, Lemme \ref{descent-lemma-descent-data-sheaves} (2),
on en déduit que $G$ est représentable, comme voulu.
\end{proof}

\begin{theorem}
\label{theorem-presentation}
Soit $S$ un schéma. Soit $U$ un schéma au-dessus de $S$.
Soit $j = (s, t) : R \to U \times_S U$
une relation d'équivalence \'etale sur $U$ au-dessus de $S$.
Alors le quotient $U/R$ est un espace algébrique,
et $U \to U/R$ est \'etale et surjectif; autrement dit,
$(U, R, U \to U/R)$ est une présentation de $U/R$.
\end{theorem}

\begin{proof}
D'après le Lemme \ref{lemma-when-it-works-it-works},
il suffit de prouver que $U/R$ est un espace algébrique.
Soit $U' \to U$ un morphisme surjectif et \'etale.
Alors $\{U' \to U\}$ est en particulier un recouvrement fppf.
Soit $R'$ la restriction de $R$ à $U'$; voir
Groupoïdes, Définition \ref{groupoids-definition-restrict-relation}.
D'après
Groupoïdes, Lemme \ref{groupoids-lemma-quotient-groupoid-restrict},
on a $U/R \cong U'/R'$.
D'après le Lemme \ref{lemma-pullback-etale-equivalence-relation}, $R'$ est une
relation d'équivalence \'etale sur $U'$. Nous pouvons donc remplacer $U$ par $U'$.

\medskip\noindent
Appliquons la remarque précédente à $U' = \coprod U_i$, où
$U = \bigcup U_i$ est un recouvrement ouvert affine de $U$. Nous
pouvons donc supposer, et nous le ferons, que $U = \coprod U_i$, où
chaque $U_i$ est un schéma affine.

\medskip\noindent
Considérons la restriction $R_i$ de $R$ à $U_i$.
D'après le Lemme \ref{lemma-pullback-etale-equivalence-relation},
c'est une relation d'équivalence \'etale.
Posons $F_i = U_i/R_i$ et $F = U/R$.
Il est clair que $\coprod F_i \to F$ est surjectif.
D'après le Lemme \ref{lemma-finding-opens}, chaque $F_i \to F$
est représentable et est une immersion ouverte.
En appliquant le Lemme \ref{lemma-presentation-quasi-compact}
à $(U_i, R_i)$, on voit que $F_i$ est un espace algébrique.
Puis, d'après le Lemme \ref{lemma-when-it-works-it-works},
$U_i \to F_i$ est \'etale et surjectif.
Il résulte du Lemme \ref{lemma-coproduct-algebraic-spaces}
que $\coprod F_i$ est un espace algébrique.
Enfin, nous avons vérifié toutes les
hypothèses du Lemme \ref{lemma-glueing-algebraic-spaces},
et il s'ensuit que $F = U/R$ est un espace algébrique.
\end{proof}










\section{Espaces algébriques, reformulation}
\label{section-algebraic-spaces-retrofitted}

\noindent
Nous commençons à constituer notre arsenal de lemmes sur les espaces algébriques.
Le premier résultat affirme que, dans la Définition \ref{definition-algebraic-space},
on peut affaiblir comme suit la condition portant sur la diagonale.

\begin{lemma}
\label{lemma-etale-locally-representable-gives-space}
Soit $S$ un schéma appartenant à $\Sch_{fppf}$.
Soit $F$ un faisceau sur $(\Sch/S)_{fppf}$
tel qu'il existe $U \in \Ob((\Sch/S)_{fppf})$ et un morphisme
$U \to F$ représentable, surjectif et \'etale.
Alors $F$ est un espace algébrique.
\end{lemma}

\begin{proof}
Posons $R = U \times_F U$. C'est un schéma, car $U \to F$ est supposé représentable.
Les projections $s, t : R \to U$ sont \'etales, car $U \to F$ est supposé \'etale.
Le morphisme $j = (t, s) : R \to U \times_S U$ est un monomorphisme et définit une relation
d'équivalence, puisque $R = U \times_F U$. D'après le Théorème \ref{theorem-presentation},
le faisceau quotient $F' = U/R$ est un espace algébrique et $U \to F'$
est surjectif et \'etale. À nouveau, comme $R = U \times_F U$, on obtient
une factorisation canonique $U \to F' \to F$, et $F' \to F$ est un morphisme injectif
de faisceaux. D'autre part, $U \to F$ est surjectif comme morphisme
de faisceaux d'après le Lemme \ref{lemma-surjective-flat-locally-finite-presentation}.
Ainsi, $F' \to F$ est également surjectif, et l'on conclut que $F' = F$ est un
espace algébrique.
\end{proof}

\begin{lemma}
\label{lemma-etale-locally-representable-by-space-gives-space}
Soit $S$ un schéma appartenant à $\Sch_{fppf}$. Soit $G$ un espace
algébrique au-dessus de $S$, soit $F$ un faisceau sur $(\Sch/S)_{fppf}$, et soit
$G \to F$ un morphisme représentable de foncteurs qui est
surjectif et \'etale. Alors $F$ est un espace algébrique.
\end{lemma}

\begin{proof}
Choisissons un schéma $U$ et un morphisme surjectif \'etale $U \to G$.
Puisque $G$ est un espace algébrique, $U \to G$ est représentable.
Le composé $U \to G \to F$ est donc représentable,
surjectif et \'etale. Voir les Lemmes
\ref{lemma-composition-representable-transformations} et
\ref{lemma-composition-representable-transformations-property}.
Ainsi, $F$ est un espace algébrique d'après le
Lemme \ref{lemma-etale-locally-representable-gives-space}.
\end{proof}

\begin{lemma}
\label{lemma-representable-over-space}
\begin{slogan}
Un foncteur qui admet un morphisme représentable vers un espace algébrique est
un espace algébrique.
\end{slogan}
Soit $S$ un schéma appartenant à $\Sch_{fppf}$.
Soit $F$ un espace algébrique au-dessus de $S$.
Soit $G \to F$ un morphisme représentable de foncteurs.
Alors $G$ est un espace algébrique.
\end{lemma}

\begin{proof}
D'après le Lemme \ref{lemma-representable-transformation-to-sheaf},
on voit que $G$ est un faisceau. Le diagramme
$$
\xymatrix{
G \times_F G \ar[r] \ar[d] & F \ar[d]^{\Delta_F} \\
G \times G \ar[r] & F \times F
}
$$
est cartésien. On voit donc que $G \times_F G \to G \times G$
est représentable d'après le
Lemme \ref{lemma-base-change-representable-transformations}.
D'après le
Lemme \ref{lemma-representable-transformation-diagonal},
on voit que $G \to G \times_F G$ est représentable.
Ainsi, $\Delta_G : G \to G \times G$ est représentable comme composé
de morphismes représentables de foncteurs; voir le
Lemme \ref{lemma-composition-representable-transformations}.
Enfin, soit $U$ un objet de $(\Sch/S)_{fppf}$
et soit $U \to F$ surjectif et \'etale. Par hypothèse,
$U \times_F G$ est représentable par un schéma $U'$. D'après le
Lemme \ref{lemma-base-change-representable-transformations-property},
le morphisme $U' \to G$ est surjectif et \'etale. Cela vérifie
la dernière condition de la Définition \ref{definition-algebraic-space}, ce qui conclut.
\end{proof}

\begin{lemma}
\label{lemma-representable-morphisms-spaces-property}
Soit $S$ un schéma appartenant à $\Sch_{fppf}$.
Soient $F$, $G$ des espaces algébriques au-dessus de $S$.
Soit $G \to F$ un morphisme représentable.
Soient $U \in \Ob((\Sch/S)_{fppf})$ et $q : U \to F$
surjectif et \'etale. Posons $V = G \times_F U$.
Enfin, soit $\mathcal{P}$ une propriété des morphismes
de schémas comme dans la Définition \ref{definition-relative-representable-property}.
Alors $G \to F$ possède la propriété $\mathcal{P}$ si et seulement si
$V \to U$ possède la propriété $\mathcal{P}$.
\end{lemma}

\begin{proof}
(Ce lemme résulte des
Lemmes \ref{lemma-base-change-representable-transformations-property} et
\ref{lemma-descent-representable-transformations-property},
mais nous en donnons aussi ici une démonstration directe.)
Il résulte immédiatement des définitions que, si $G \to F$ possède la propriété
$\mathcal{P}$, alors $V \to U$ possède la propriété $\mathcal{P}$.
Réciproquement, supposons que $V \to U$ possède la propriété $\mathcal{P}$.
Soit $T \to F$ un morphisme d'un schéma vers $F$.
Soit $T' = T \times_F G$, qui est un schéma puisque $G \to F$ est
représentable. Il faut montrer que $T' \to T$ possède la propriété $\mathcal{P}$.
Considérons le diagramme commutatif de schémas
$$
\xymatrix{
V \ar[d] & T \times_F V \ar[d] \ar[l] \ar[r] &
T \times_F G \ar[d] \ar@{=}[r] & T' \\
U & T \times_F U \ar[l] \ar[r] & T
}
$$
où les deux carrés sont cartésiens. On en déduit que
la flèche médiane possède la propriété $\mathcal{P}$ en tant que changement de base
de $V \to U$. Enfin, $\{T \times_F U \to T\}$ est un recouvrement fppf,
car il est surjectif \'etale; on en déduit donc que
$T' \to T$ possède la propriété $\mathcal{P}$, puisque celle-ci est locale sur la
base pour la topologie fppf.
\end{proof}

\begin{lemma}
\label{lemma-morphism-sheaves-with-P-effective-descent-etale}
Soit $S$ un schéma appartenant à $\Sch_{fppf}$.
Soit $G \to F$ un morphisme de préfaisceaux sur $(\Sch/S)_{fppf}$.
Soit $\mathcal{P}$ une propriété des morphismes de schémas.
Supposons que
\begin{enumerate}
\item $\mathcal{P}$ soit stable par tout changement de base, locale sur la
base pour la topologie fppf, et que les morphismes de type $\mathcal{P}$ satisfassent à la descente pour les recouvrements fppf;
voir Descente, Définition \ref{descent-definition-descending-types-morphisms},
\item $G$ soit un faisceau,
\item $F$ soit un espace algébrique,
\item il existe $U \in \Ob((\Sch/S)_{fppf})$
et un morphisme surjectif \'etale $U \to F$ tels que
$V = G \times_F U$ soit représentable, et
\item $V \to U$ possède la propriété $\mathcal{P}$.
\end{enumerate}
Alors $G$ est un espace algébrique, $G \to F$ est représentable et possède la propriété
$\mathcal{P}$.
\end{lemma}

\begin{proof}
Soit $T$ un schéma et soit $T \to F$ un morphisme. Alors
$U \times_F T \to T$ est surjectif \'etale; ainsi, $\{U \times_F T \to T\}$ est
un recouvrement pour la topologie \'etale. Considérons
$$
W = G \times_F (U \times_F T) = V \times_F T = V \times_U (U \times_F T).
$$
C'est un schéma puisque $F$ est un espace algébrique. Le morphisme
$W \to U \times_F T$ possède la propriété $\mathcal{P}$, car c'est un
changement de base de $V \to U$. Il existe un isomorphisme
\begin{align*}
W \times_T (U \times_F T) & =
(G \times_F (U \times_F T)) \times_T (U \times_F T) \\
& = (U \times_F T) \times_T (G \times_F (U \times_F T)) \\
& = (U \times_F T) \times_T W
\end{align*}
au-dessus de $(U \times_F T) \times_T (U \times_F T)$. L'égalité médiane envoie
$((g, (u_1, t)), (u_2, t))$ sur $((u_1, t), (g, (u_2, t)))$.
Cela définit une donnée de descente pour $W/U \times_F T/T$; voir
Descente, Définition \ref{descent-definition-descent-datum}.
Cela résulte du
Lemme \ref{descent-lemma-descent-data-sheaves} de Descente.
Plus précisément, on dispose du faisceau $G \times_F T$, dont le
changement de base à $U \times_F T$ est représenté par $W$, et l'isomorphisme
ci-dessus est celui de la démonstration du
Lemme \ref{descent-lemma-descent-data-sheaves} de Descente.
Par hypothèse sur $\mathcal{P}$, la donnée de descente ci-dessus est représentable.
La dernière assertion du
Lemme \ref{descent-lemma-descent-data-sheaves} de Descente
montre donc que $G \times_F T$ est représentable. Cela prouve que
$G \to F$ est un morphisme représentable de foncteurs.

\medskip\noindent
Comme $G \to F$ est représentable, on voit que $G$ est un espace algébrique d'après le
Lemme \ref{lemma-representable-over-space}. Le fait que $G \to F$ possède
la propriété $\mathcal{P}$ résulte maintenant du
Lemme \ref{lemma-representable-morphisms-spaces-property}.
\end{proof}

\begin{lemma}
\label{lemma-lift-morphism-presentations}
Soit $S$ un schéma appartenant à $\Sch_{fppf}$.
Soient $F, G$ des espaces algébriques au-dessus de $S$.
Soit $a : F \to G$ un morphisme.
Étant donnés $V \in \Ob((\Sch/S)_{fppf})$
et un morphisme surjectif \'etale $q : V \to G$, il existe
$U \in \Ob((\Sch/S)_{fppf})$
et un diagramme commutatif
$$
\xymatrix{
U \ar[d]_p \ar[r]_\alpha &
V \ar[d]^q \\
F \ar[r]^a & G
}
$$
où $p$ est surjectif et \'etale.
\end{lemma}

\begin{proof}
Choisissons d'abord $W \in \Ob((\Sch/S)_{fppf})$
et un morphisme surjectif \'etale $W \to F$.
Posons ensuite $U = W \times_G V$. Puisque $G$ est un espace algébrique,
on voit que $U$ est isomorphe à un objet de $(\Sch/S)_{fppf}$.
Comme $q$ est surjectif \'etale, on voit que $U \to W$ est surjectif
\'etale (voir le
Lemme \ref{lemma-base-change-representable-transformations-property}).
Ainsi, $U \to F$ est surjectif \'etale comme composé de morphismes surjectifs
\'etales (voir le
Lemme \ref{lemma-composition-representable-transformations-property}).
\end{proof}



























\section{Immersions et recouvrements de Zariski des espaces algébriques}
\sectionmark{Immersions et recouvrements de Zariski}
\label{section-Zariski}

\noindent
À ce stade, un phénomène intéressant se produit. Nous avons déjà défini
la notion d'immersion ouverte d'espaces algébriques (au moyen de la
Définition \ref{definition-relative-representable-property}),
mais il nous reste à définir la notion de {\it point}\footnote{Nous
associerons un espace topologique à un espace algébrique dans
Propriétés des espaces algébriques, section
\ref{spaces-properties-section-points},
et ses ouverts correspondront exactement aux sous-espaces ouverts définis ci-dessous.}.
Ainsi, la {\it topologie de Zariski} d'un espace algébrique
a déjà été définie, alors qu'il n'y a encore aucun espace !

\medskip\noindent
Peut-être de façon superflue, nous introduisons formellement les immersions comme suit.

\begin{definition}
\label{definition-immersion}
Soit $S \in \Ob(\Sch_{fppf})$ un schéma.
Soit $F$ un espace algébrique au-dessus de $S$.
\begin{enumerate}
\item Un morphisme d'espaces algébriques au-dessus de $S$
est appelé {\it immersion ouverte} s'il est représentable et si c'est une immersion ouverte
au sens de la Définition \ref{definition-relative-representable-property}.
\item Un {\it sous-espace ouvert} de $F$ est un sous-foncteur $F' \subset F$
tel que $F'$ soit un espace algébrique et que $F' \to F$ soit une
immersion ouverte.
\item Un morphisme d'espaces algébriques au-dessus de $S$
est appelé {\it immersion fermée} s'il est représentable et si c'est une immersion
fermée au sens de la
Définition \ref{definition-relative-representable-property}.
\item Un {\it sous-espace fermé} de $F$ est un sous-foncteur $F' \subset F$
tel que $F'$ soit un espace algébrique et que $F' \to F$ soit une
immersion fermée.
\item Un morphisme d'espaces algébriques au-dessus de $S$
est appelé {\it immersion} s'il est représentable et si c'est une immersion
au sens de la Définition \ref{definition-relative-representable-property}.
\item Un {\it sous-espace localement fermé} de $F$ est un sous-foncteur $F' \subset F$
tel que $F'$ soit un espace algébrique et que $F' \to F$ soit une
immersion.
\end{enumerate}
\end{definition}

\noindent
Notons que ces définitions ont un sens, puisqu'une immersion
est en particulier un monomorphisme (voir
Schémas, Lemme \ref{schemes-lemma-immersions-monomorphisms}
et Lemme \ref{lemma-representable-transformations-property-implication}),
et que, par conséquent, l'image d'une
immersion $G \to F$ d'espaces algébriques est un sous-foncteur $F' \subset F$
qui est (canoniquement) isomorphe à $G$. Ainsi, une partie de la discussion
de Schémas, section \ref{schemes-section-immersions}, se transpose au
cadre des espaces algébriques.

\begin{lemma}
\label{lemma-composition-immersions}
Soit $S \in \Ob(\Sch_{fppf})$ un schéma.
Un composé d'immersions (fermées, resp.\ ouvertes)
d'espaces algébriques au-dessus de $S$ est une immersion (fermée, resp.\ ouverte)
d'espaces algébriques au-dessus de $S$.
\end{lemma}

\begin{proof}
Voir le Lemme \ref{lemma-composition-representable-transformations-property} et les
Remarques \ref{remark-list-properties-fpqc-local-base} (voir la toute dernière ligne de
cette remarque) et \ref{remark-list-properties-stable-composition}.
\end{proof}

\begin{lemma}
\label{lemma-base-change-immersions}
Soit $S \in \Ob(\Sch_{fppf})$ un schéma.
Un changement de base d'une immersion (fermée, resp.\ ouverte)
d'espaces algébriques au-dessus de $S$ est une immersion (fermée, resp.\ ouverte)
d'espaces algébriques au-dessus de $S$.
\end{lemma}

\begin{proof}
Voir le Lemme \ref{lemma-base-change-representable-transformations-property} et la
Remarque \ref{remark-list-properties-fpqc-local-base} (voir la toute dernière ligne de
cette remarque).
\end{proof}

\begin{lemma}
\label{lemma-sub-subspaces}
Soit $S \in \Ob(\Sch_{fppf})$ un schéma.
Soit $F$ un espace algébrique au-dessus de $S$. Soient $F_1$, $F_2$ des
sous-espaces localement fermés de $F$. Si $F_1 \subset F_2$ comme sous-foncteurs
de $F$, alors $F_1$ est un sous-espace localement fermé de $F_2$.
Il en va de même pour les sous-espaces fermés et ouverts.
\end{lemma}

\begin{proof}
Soit $T \to F_2$ un morphisme, où $T$ est un schéma.
Puisque $F_2 \to F$ est un monomorphisme, on voit que
$T \times_{F_2} F_1 = T \times_F F_1$. Le lemme en résulte
formellement.
\end{proof}

\noindent
Définissons formellement la notion de recouvrement ouvert de Zariski
d'espaces algébriques. Notons que, dans le Lemme \ref{lemma-glueing-algebraic-spaces},
nous avons déjà rencontré de tels recouvrements ouverts comme méthode pour
construire des espaces algébriques.

\begin{definition}
\label{definition-Zariski-open-covering}
Soit $S \in \Ob(\Sch_{fppf})$ un schéma.
Soit $F$ un espace algébrique au-dessus de $S$.
Un {\it recouvrement de Zariski} $\{F_i \subset F\}_{i \in I}$ de $F$
est constitué d'un ensemble $I$ et d'une famille de sous-espaces ouverts
$F_i \subset F$ telle que $\coprod F_i \to F$ soit un morphisme surjectif
de faisceaux.
\end{definition}

\noindent
Notons que, si $T$ est un schéma
et si $a : T \to F$ est un morphisme, chacun des produits fibrés
$T \times_F F_i$ s'identifie à un sous-schéma ouvert
$T_i \subset T$. La dernière condition de la définition signifie
exactement que $T = \bigcup_{i \in I} T_i$.

\medskip\noindent
Il est clair que la collection $F_{Zar}$ des sous-espaces ouverts de
$F$ est un ensemble (puisque $(\Sch/S)_{fppf}$ est un site, donc un ensemble).
De plus, on peut faire de $F_{Zar}$ une catégorie en prenant pour
morphismes les inclusions de sous-foncteurs (qui sont automatiquement des immersions
ouvertes d'après le Lemme \ref{lemma-sub-subspaces}). Enfin, la
Définition \ref{definition-Zariski-open-covering}
fournit la notion de recouvrement de Zariski $\{F_i \to F'\}_{i \in I}$
dans la catégorie $F_{Zar}$. Ainsi, comme dans le cas d'un espace topologique
(voir Sites, Exemple \ref{sites-example-site-topological}),
en choisissant convenablement un ensemble de recouvrements,
on peut obtenir un site de Zariski de l'espace algébrique $F$.

\begin{definition}
\label{definition-small-Zariski-site}
Soit $S \in \Ob(\Sch_{fppf})$ un schéma. Soit $F$ un espace algébrique au-dessus de
$S$. Un {\it petit site de Zariski $F_{Zar}$} d'un espace algébrique $F$ est l'un
des sites décrits ci-dessus.
\end{definition}

\noindent
Cela donne donc un sens à l'expression : une propriété est vraie
localement pour la topologie de Zariski sur un espace algébrique; c'est ainsi que nous utiliserons cette
notion. En général, la topologie de Zariski n'est pas assez fine pour nos
besoins. Par exemple, on peut considérer la catégorie des faisceaux de Zariski
sur un espace algébrique. Il s'avérera que ce n'est pas la
bonne notion à considérer, même pour les faisceaux quasi-cohérents.
On n'obtient le résultat recherché qu'en utilisant le site \'etale ou fppf de
$F$ pour définir les faisceaux quasi-cohérents.











\section{Conditions de séparation des espaces algébriques}
\label{section-separation}

\noindent
Une condition de séparation pour un espace algébrique $F$ est une condition
portant sur le morphisme diagonal $F \to F \times F$. Commençons par
énumérer les propriétés que la diagonale possède automatiquement.
Comme la diagonale est représentable par définition, l'énoncé du lemme suivant
a un sens (grâce à la
Définition \ref{definition-relative-representable-property}).

\begin{lemma}
\label{lemma-properties-diagonal}
Soit $S$ un schéma appartenant à $\Sch_{fppf}$.
Soit $F$ un espace algébrique sur $S$.
Soit $\Delta : F \to F \times F$ le morphisme diagonal.
Alors
\begin{enumerate}
\item $\Delta$ est localement de type fini,
\item $\Delta$ est un monomorphisme,
\item $\Delta$ est séparé, et
\item $\Delta$ est localement quasi-fini.
\end{enumerate}
\end{lemma}

\begin{proof}
Soit $F = U/R$ une présentation de $F$.
Comme dans la démonstration du Lemme \ref{lemma-presentation-quasi-compact}, le diagramme
$$
\xymatrix{
R \ar[r] \ar[d]_j & F \ar[d]^\Delta \\
U \times_S U \ar[r] & F \times F
}
$$
est cartésien. Par conséquent, d'après le
Lemme \ref{lemma-representable-morphisms-spaces-property},
il suffit de montrer que $j$ possède les propriétés énumérées dans le lemme.
(Notons que chacune des propriétés (1) -- (4) figure dans les listes
des Remarques \ref{remark-list-properties-stable-base-change}
et \ref{remark-list-properties-fpqc-local-base}.)
Comme $j$ définit une relation d'équivalence, c'est un monomorphisme.
Il est donc séparé d'après
Schémas, Lemme \ref{schemes-lemma-monomorphism-separated}.
Comme $R$ est une relation d'équivalence \'etale, on voit que
$s, t : R \to U$ sont \'etales. Ainsi, $s, t$ sont localement de
type fini. Il résulte alors de
Morphismes, Lemme \ref{morphisms-lemma-permanence-finite-type}, que
$j$ est localement de type fini. Enfin, comme c'est un monomorphisme,
ses fibres sont finies. On en conclut qu'il est localement quasi-fini d'après
Morphismes, Lemme \ref{morphisms-lemma-finite-fibre}.
\end{proof}

\noindent
Voici quelques types usuels de conditions de séparation, relativement au schéma
de base $S$. Il existe aussi une notion absolue de ces conditions, dont nous
discuterons dans
Propriétés des espaces algébriques, section \ref{spaces-properties-section-separation}.
En outre, nous étudierons les conditions de séparation d'un morphisme
d'espaces algébriques dans
Morphismes d'espaces algébriques, section \ref{spaces-morphisms-section-separation-axioms}.

\begin{definition}
\label{definition-separated}
Soit $S$ un schéma appartenant à $\Sch_{fppf}$.
Soit $F$ un espace algébrique sur $S$.
Soit $\Delta : F \to F \times F$ le morphisme diagonal.
\begin{enumerate}
\item On dit que $F$ est {\it séparé sur $S$} si $\Delta$ est une immersion fermée.
\item On dit que $F$ est {\it localement séparé sur $S$}\footnote{Dans la
littérature, cela désigne souvent des espaces algébriques quasi-séparés et
localement séparés.} si $\Delta$ est une
immersion.
\item On dit que $F$ est {\it quasi-séparé sur $S$} si $\Delta$ est quasi-compact.
\item On dit que $F$ est {\it localement quasi-séparé sur $S$ pour la topologie de Zariski}\footnote{Cette
définition a été proposée par B.\ Conrad.} s'il
existe un recouvrement de Zariski $F = \bigcup_{i \in I} F_i$ tel que
chacun des $F_i$ soit quasi-séparé.
\end{enumerate}
\end{definition}

\noindent
Notons que, si la diagonale est quasi-compacte (lorsque $F$ est séparé ou
quasi-séparé), elle est en fait
quasi-finie et séparée, donc quasi-affine (d'après Compléments sur les morphismes,
Lemme \ref{more-morphisms-lemma-quasi-finite-separated-quasi-affine}).








\section{Exemples d'espaces algébriques}
\label{section-examples}

\noindent
Dans cette section, nous construisons quelques exemples d'espaces algébriques.
Certains de ces exemples ont été suggérés par B.\ Conrad.
Comme nous ne disposons pas encore de beaucoup de théorie,
l'exposé est quelque peu malaisé par endroits.

\begin{example}
\label{example-affine-line-involution}
Soit $k$ un corps de caractéristique $\not = 2$. Soit $U = \mathbf{A}^1_k$. Posons
$$
j : R = \Delta \amalg \Gamma \longrightarrow U \times_k U
$$
où $\Delta = \{(x, x) \mid x \in \mathbf{A}^1_k\}$ et
$\Gamma = \{(x, -x) \mid x \in \mathbf{A}^1_k, x \not = 0\}$.
Il est clair que $s, t : R \to U$ sont \'etales, et donc que
$j$ est une relation d'équivalence \'etale.
Le quotient $X = U/R$ est un espace algébrique d'après le
Théorème \ref{theorem-presentation}.
Comme $R$ est quasi-compact, on voit que $X$ est quasi-séparé.
En revanche, $X$ n'est pas localement séparé, car
le morphisme $j$ n'est pas une immersion.
\end{example}

\begin{example}
\label{example-non-representable-descent}
Soit $k$ un corps. Soit $k'/k$ une extension galoisienne de degré $2$
avec $\text{Gal}(k'/k) = \{1, \sigma\}$. Soient $S = \Spec(k[x])$
et $U = \Spec(k'[x])$. Notons que
$$
U \times_S U =
\Spec((k' \otimes_k k')[x]) =
\Delta(U) \amalg \Delta'(U)
$$
où $\Delta' = (1, \sigma) : U \to U \times_S U$. Prenons
$$
R = \Delta(U) \amalg \Delta'(U \setminus \{0_U\})
$$
où $0_U \in U$ désigne le point $k'$-rationnel dont la coordonnée $x$ est nulle.
Il est facile de voir que $R$ est une relation d'équivalence \'etale sur $U$ au-dessus de $S$
et donc que $X = U/R$ est un espace algébrique d'après le
Théorème \ref{theorem-presentation}. Voici quelques propriétés de $X$ (dont certaines
n'auront de sens que plus tard) :
\begin{enumerate}
\item $X \to S$ est un isomorphisme au-dessus de $S \setminus \{0_S\}$,
\item le morphisme $X \to S$ est \'etale (voir
Propriétés des espaces algébriques,
Définition \ref{spaces-properties-definition-etale}),
\item la fibre $0_X$ de $X \to S$ au-dessus de $0_S$ est isomorphe à
$\Spec(k') = 0_U$,
\item $X$ n'est pas un schéma, car, s'il l'était, $\mathcal{O}_{X, 0_X}$
serait un anneau local intègre $(\mathcal{O}, \mathfrak m, \kappa)$ de
corps des fractions $k(x)$, avec $x \in \mathfrak m$, et de corps résiduel
$\kappa = k'$, ce qui est impossible,
\item $X$ n'est pas séparé, mais il est
localement séparé et quasi-séparé,
\item il existe un morphisme surjectif, fini et \'etale $S' \to S$
tel que le changement de base $X' = S' \times_S X$ soit un schéma (en effet, si
l'on effectue le changement de base vers $S' = \Spec(k'[x])$, alors $U$ se scinde en
deux copies de $S'$ et $X'$ devient isomorphe à la droite affine dont le point
$0$ est dédoublé, voir
Schémas, Exemple \ref{schemes-example-affine-space-zero-doubled}), et
\item si l'on considère $X$ comme un espace algébrique de type fini sur
$\Spec(k)$, alors, de même, le changement de base $X_{k'}$ est un schéma,
mais $X$ n'est pas un schéma.
\end{enumerate}
En particulier, on obtient ainsi un exemple de donnée de descente pour des schémas,
relativement au recouvrement $\{\Spec(k') \to \Spec(k)\}$,
qui n'est pas effective.
\end{example}

\noindent
Voir aussi Exemples,
Lemme \ref{examples-lemma-non-effective-descent-projective},
qui montre que les données de descente ne sont pas nécessairement effectives,
même pour un morphisme projectif de schémas. Cet exemple
fournit un espace algébrique lisse et séparé de dimension 3
sur ${\mathbf C}$ qui n'est pas un schéma.

\medskip\noindent
Nous utiliserons le lemme suivant comme moyen commode de construire
des espaces algébriques comme quotients de schémas par des actions libres de groupes.

\begin{lemma}
\label{lemma-quotient}
Soit $U \to S$ un morphisme de $\Sch_{fppf}$.
Soit $G$ un groupe abstrait. Soit $G \to \text{Aut}_S(U)$
un homomorphisme de groupes. Supposons que
\begin{itemize}
\item[(*)] si $u \in U$ est un point et si $g(u) = u$
pour un élément non neutre $g \in G$, alors $g$
induit un automorphisme non trivial de $\kappa(u)$.
\end{itemize}
Alors
$$
j :
R = \coprod\nolimits_{g \in G} U
\longrightarrow
U \times_S U,
\quad
(g, x) \longmapsto (g(x), x)
$$
est une relation d'équivalence \'etale et, par conséquent,
$$
F = U/R
$$
est un espace algébrique d'après le Théorème \ref{theorem-presentation}.
\end{lemma}

\begin{proof}
Dans l'énoncé du lemme, le symbole $\text{Aut}_S(U)$ désigne
le groupe des automorphismes de $U$ au-dessus de $S$.
Supposons $(*)$ vérifiée. Montrons que
$$
j :
R = \coprod\nolimits_{g \in G} U
\longrightarrow
U \times_S U,
\quad
(g, x) \longmapsto (g(x), x)
$$
est un monomorphisme. Cela signifie que, si $T$ est un schéma
non vide et si $h : T \to U$ est un point à valeurs dans $T$ tel que
$g \circ h = g' \circ h$, alors $g = g'$. Supposons
$T \not = \emptyset$, $h : T \to U$ et $g \circ h = g' \circ h$.
Soit $t \in T$. Considérons le composé
$\Spec(\kappa(t)) \to \Spec(\kappa(h(t))) \to U$.
On en conclut que $g^{-1} \circ g'$ fixe $u = h(t)$ et
agit comme l'identité sur son corps résiduel. On a donc $g = g'$ par $(*)$.

\medskip\noindent
Ainsi, si $(*)$ est vérifiée, on voit que $j$ est une relation (voir
Schémas en groupoïdes, Définition \ref{groupoids-definition-equivalence-relation}).
De plus, c'est une relation d'équivalence puisque, sur les points à valeurs dans $T$
pour un schéma connexe $T$, on a
$R(T) = G \times U(T) \to U(T) \times U(T)$ (rappelons que nous
travaillons toujours au-dessus de $S$). En outre, les morphismes $s, t : R \to U$ sont \'etales,
puisque $R$ est une somme disjointe de copies de $U$.
Cela prouve que $j : R \to U \times_S U$ est une relation d'équivalence \'etale.
\end{proof}

\noindent
Étant donnés un schéma $U$ et une action d'un groupe $G$ sur $U$, on dit que l'action
de $G$ sur $U$ est {\it libre} si la condition $(*)$ du Lemme \ref{lemma-quotient}
est vérifiée. Cela équivaut à la notion d'action libre du schéma en groupes
constant $G_S$ sur $U$, telle qu'elle est définie dans
Schémas en groupoïdes, Définition \ref{groupoids-definition-free-action}.
On peut interpréter le lemme comme affirmant que les quotients de schémas par des
actions libres de groupes existent dans la catégorie des espaces algébriques.

\begin{definition}
\label{definition-quotient}
Reprenons les notations $U \to S$, $G$, $R$ du Lemme \ref{lemma-quotient}.
Si l'action de $G$ sur $U$ satisfait $(*)$, on dit que $G$ {\it agit librement}
sur le schéma $U$. Dans ce cas, l'espace algébrique $U/R$ est noté
$U/G$ et appelé le {\it quotient de $U$ par $G$}.
\end{definition}

\noindent
Cette notation est compatible avec la notation $U/G$ introduite dans
Schémas en groupoïdes, Définition \ref{groupoids-definition-quotient-sheaf}.
Nous donnerons plus tard un sens au quotient comme champ algébrique sans
aucune hypothèse sur l'action; nous utiliserons alors la
notation $[U/G]$. Avant d'aborder les exemples, démontrons
encore quelques lemmes afin de faciliter la discussion. Le lemme suivant traite des
diverses conditions de séparation de ce quotient lorsque $G$ est fini.

\begin{lemma}
\label{lemma-quotient-finite-separated}
Reprenons les notations et les hypothèses du Lemme \ref{lemma-quotient}.
Supposons $G$ fini. Alors
\begin{enumerate}
\item si $U \to S$ est quasi-séparé, alors $U/G$ est quasi-séparé
sur $S$, et
\item si $U \to S$ est séparé, alors $U/G$ est séparé sur $S$.
\end{enumerate}
\end{lemma}

\begin{proof}
Dans la démonstration du Lemme \ref{lemma-properties-diagonal},
nous avons vu qu'il suffit d'établir les
propriétés correspondantes pour le morphisme $j : R \to U \times_S U$.
Si $U \to S$ est quasi-séparé, alors, pour tout ouvert affine $V \subset U$
dont l'image est contenue dans un ouvert affine de $S$,
les ouverts $g(V) \cap V$ sont quasi-compacts. Il s'ensuit que $j$ est
quasi-compact.
Si $U \to S$ est séparé, la diagonale $\Delta_{U/S}$ est une immersion
fermée. Ainsi, $j : R \to U \times_S U$ est une somme finie
d'immersions fermées aux images disjointes. Ainsi, $j$ est une immersion fermée.
\end{proof}

\begin{lemma}
\label{lemma-quotient-field-map}
Reprenons les notations et les hypothèses du Lemme \ref{lemma-quotient}.
Si $\Spec(k) \to U/G$ est un morphisme, alors il existe
\begin{enumerate}
\item une extension galoisienne finie $k'/k$,
\item un sous-groupe fini $H \subset G$,
\item un isomorphisme $H \to \text{Gal}(k'/k)$, et
\item un morphisme $H$-équivariant $\Spec(k') \to U$.
\end{enumerate}
Réciproquement, de telles données déterminent un morphisme $\Spec(k) \to U/G$.
\end{lemma}

\begin{proof}
Considérons le produit fibré $V = \Spec(k) \times_{U/G} U$.
Voici un diagramme
$$
\xymatrix{
V \ar[r] \ar[d] & U \ar[d] \\
\Spec(k) \ar[r] & U/G
}
$$
Alors $V$ est un schéma non vide \'etale sur $\Spec(k)$ et est donc une
réunion disjointe $V = \coprod_{i \in I} \Spec(k_i)$
de spectres de corps $k_i$, extensions finies séparables de $k$
(Morphismes de schémas, Lemme \ref{morphisms-lemma-etale-over-field}).
On a
\begin{align*}
V \times_{\Spec(k)} V
& =
(\Spec(k) \times_{U/G} U) \times_{\Spec(k)}(\Spec(k) \times_{U/G} U) \\
& = 
\Spec(k) \times_{U/G} U \times_{U/G} U \\
& =
\Spec(k) \times_{U/G} U \times G \\
& =
V \times G
\end{align*}
L'action de $G$ sur $U$ induit une action $a : G \times V \to V$.
L'égalité affichée signifie que
$G \times V \to V \times_{\Spec(k)} V$, $(g, v) \mapsto (a(g, v), v)$
est un isomorphisme. En particulier, on voit que, pour tout $i$, on a
un isomorphisme $H_i \times \Spec(k_i) \to \Spec(k_i \otimes_k k_i)$,
où $H_i \subset G$ est le sous-groupe des éléments qui fixent $i \in I$.
Ainsi, $H_i$ est fini et est le groupe de Galois de $k_i/k$.
Nous omettons la construction réciproque.
\end{proof}

\noindent
Il résulte par exemple de ce lemme que,
si $k'/k$ est une extension galoisienne finie, alors
$\Spec(k')/\text{Gal}(k'/k) \cong \Spec(k)$.
Que se passe-t-il si l'extension est infinie ? Voici un exemple.

\begin{example}
\label{example-Qbar}
Soit $S = \Spec(\mathbf{Q})$.
Soit $U = \Spec(\overline{\mathbf{Q}})$.
Soit $G = \text{Gal}(\overline{\mathbf{Q}}/\mathbf{Q})$, muni de son action
évidente sur $U$. Alors, par construction, la propriété $(*)$ du
Lemme \ref{lemma-quotient} est vérifiée, et l'on obtient un espace algébrique
$$
X = \Spec(\overline{\mathbf{Q}})/G
\longrightarrow
S = \Spec(\mathbf{Q}).
$$
Bien entendu, c'est une approximation tout à fait ridicule de $S$ !
En effet, d'après le théorème d'Artin-Schreier,
voir \cite[Théorème 17, page 316]{JacobsonIII},
les seuls sous-groupes finis de $\text{Gal}(\overline{\mathbf{Q}}/\mathbf{Q})$
sont $\{1\}$ et les conjugués du groupe d'ordre deux
$\text{Gal}(\overline{\mathbf{Q}}/\overline{\mathbf{Q}} \cap \mathbf{R})$.
Par conséquent, si
$\Spec(k) \to X$ est un morphisme avec $k$ algébrique sur $\mathbf{Q}$,
il résulte du Lemme \ref{lemma-quotient-field-map} et du théorème
que nous venons de mentionner que $k$ est soit $\overline{\mathbf{Q}}$, soit isomorphe à
$\overline{\mathbf{Q}} \cap \mathbf{R}$.
\end{example}

\noindent
Ce qui ne va pas dans l'exemple ci-dessus, c'est que
le groupe de Galois est muni d'une topologie,
et que celle-ci devrait intervenir d'une manière ou d'une autre dans toute construction
d'un quotient de $\Spec(\overline{\mathbf{Q}})$.
L'exemple suivant est, à mon avis, beaucoup plus raisonnable
et peut effectivement apparaître dans « la nature ».

\begin{example}
\label{example-affine-line-translation}
Soit $k$ un corps de caractéristique zéro.
Soit $U = \mathbf{A}^1_k$ et soit $G = \mathbf{Z}$.
Nous prenons pour action $n(x) = x + n$, c'est-à-dire l'action de
$\mathbf{Z}$ sur la droite affine par translation.
Le seul point fixe est le point générique, et il
est clair que $\mathbf{Z}$ s'injecte dans
le groupe des automorphismes du corps $k(x)$. (C'est
ici que nous utilisons l'hypothèse de caractéristique zéro.)
Considérons le morphisme
$$
\gamma : \Spec(k(x)) \longrightarrow X = \mathbf{A}^1_k/\mathbf{Z}
$$
du point générique de la droite affine vers le quotient.
Nous affirmons que ce morphisme ne se factorise par aucun
monomorphisme $\Spec(L) \to X$ du spectre d'un
corps vers $X$. (À la différence de ce qui se passe pour les schémas, voir
Schémas, section \ref{schemes-section-points}.) En effet, comme
$\mathbf{Z}$ ne possède aucun sous-groupe fini non trivial, il résulte du
Lemme \ref{lemma-quotient-field-map} que, pour toute telle
factorisation, $k(x) = L$. Enfin, $\gamma$ n'est pas un monomorphisme,
puisque
$$
\Spec(k(x)) \times_{\gamma, X, \gamma} \Spec(k(x))
\cong
\Spec(k(x)) \times \mathbf{Z}.
$$
\end{example}

\noindent
Cet exemple suggère que, pour définir les points d'un espace algébrique
$X$, il faut considérer les classes d'équivalence de morphismes de spectres
de corps vers $X$, et non l'ensemble des monomorphismes de spectres de corps.

\medskip\noindent
Nous terminons par un exemple véritablement affreux.

\begin{example}
\label{example-infinite-product}
Soit $k$ un corps.
Soit $A = \prod_{n \in \mathbf{N}} k$ le produit infini.
Posons $U = \Spec(A)$, considéré comme un schéma sur $S = \Spec(k)$.
Notons que les projections $\text{pr}_n : A \to k$ définissent des immersions
ouvertes et fermées $f_n : S \to U$. Posons
$$
R = U \amalg \coprod\nolimits_{(n, m) \in \mathbf{N}^2, \ n \not = m} S
$$
où le morphisme $j$ est égal à $\Delta_{U/S}$ sur la composante $U$
et à $j = (f_n, f_m)$ sur la composante $S$ correspondant à $(n, m)$.
La remarque ci-dessus montre clairement que $s, t$ sont \'etales.
Il est également clair que $j$ est une relation d'équivalence. Nous
obtenons donc un espace algébrique
$$
X = U/R.
$$
Pour comprendre ce que cela signifie, plaçons-nous dans le cas où
le corps $k$ est fini et possède $q$ éléments. Commençons par
examiner quelque peu l'espace topologique $|U|$ associé au schéma $U$.
Tout élément de $A$ vérifie $x^q = x$.
Par conséquent, tout corps résiduel de $A$ est isomorphe à $k$, et
tous les points de $U$ sont fermés. Mais la topologie de $U$ n'est pas
la topologie discrète. Soit $u_n \in |U|$ le point correspondant
à $f_n$. Comme indiqué ci-dessus, les points $u_n$ sont
les points ouverts (et sont donc isolés). Il doit donc exister
d'autres points, puisque nous savons que $U$ est quasi-compact, voir
Algèbre commutative, Lemme \ref{algebra-lemma-quasi-compact}
(et n'est donc pas égal à un ensemble discret infini).
Une autre façon de le voir consiste à remarquer que l'idéal (propre)
$$
I =
\{x = (x_n) \in A \mid \text{tous les }x_n\text{ sauf un nombre fini d'entre eux sont nuls}\}
$$
est contenu dans un idéal maximal. Notons aussi que tout élément
$x$ de $A$ s'écrit $x = ue$, où $u$ est une unité et $e$ un
idempotent. Une base de la topologie de $A$ est donc formée de parties ouvertes et
fermées (voir Algèbre commutative, Lemme \ref{algebra-lemma-idempotent-spec}.)
Ainsi, $|U|$ est totalement discontinu, mais non trivial.
Enfin, notons que $\{u_n\}$ est dense dans $|U|$.

\medskip\noindent
Nous définirons plus tard un espace topologique $|X|$ associé à $X$, voir
Propriétés des espaces algébriques, section \ref{spaces-properties-section-points}.
Que peut-on dire de $|X|$ ?
Il s'avère que l'application $|U| \to |X|$ est surjective et continue.
Tous les points $u_n$ ont pour image le même point $x_0$ de $|X|$, et aucun
des autres points n'est identifié à un autre. Comme $\{u_n\}$ est dense dans $|U|$, on
conclut que l'adhérence de $x_0$ dans $|X|$ est $|X|$. Autrement dit,
$|X|$ est irréductible et $x_0$ est un point générique de $|X|$. Cela paraît
bizarre, puisque $x_0$ est aussi l'image d'une section
$S \to X$ du morphisme structural $X \to S$ (et, dans le cas des
schémas, cela impliquerait qu'il s'agit d'un point fermé, voir
Morphismes de schémas, Lemme
\ref{morphisms-lemma-algebraic-residue-field-extension-closed-point-fibre}).
\end{example}

\noindent
Quoi que l'on pense de ce qui se passe réellement dans cet exemple, il montre assurément
qu'il faut faire preuve d'une certaine prudence lorsque l'on définit les composantes
irréductibles, la connexité, etc., des espaces algébriques.






\section{Changement de grands sites}
\label{section-change-big-site}

\noindent
Dans cette section, nous examinons brièvement ce qui se passe lorsque l'on
change de grand site. En substance, on peut toujours agrandir à volonté le grand
site; on peut donc supposer que tout ensemble de schémas que l'on souhaite
considérer est contenu dans le grand site fppf sur lequel on considère notre
espace algébrique. Voici un énoncé précis du résultat.

\begin{lemma}
\label{lemma-change-big-site}
Supposons donnés de grands sites $\Sch_{fppf}$ et $\Sch'_{fppf}$.
Supposons que $\Sch_{fppf}$ soit contenu dans $\Sch'_{fppf}$,
voir Topologies sur les schémas, section \ref{topologies-section-change-alpha}.
Soit $S$ un objet de $\Sch_{fppf}$. Soient
\begin{align*}
g : \Sh((\Sch/S)_{fppf})
\longrightarrow
\Sh((\Sch'/S)_{fppf}), \\
f : \Sh((\Sch'/S)_{fppf})
\longrightarrow
\Sh((\Sch/S)_{fppf})
\end{align*}
les morphismes de topos de
Topologies sur les schémas, Lemme \ref{topologies-lemma-change-alpha}.
Soit $F$ un faisceau d'ensembles sur $(\Sch/S)_{fppf}$. Alors
\begin{enumerate}
\item si $F$ est représentable par un schéma
$X \in \Ob((\Sch/S)_{fppf})$ sur $S$,
alors $f^{-1}F$ est lui aussi représentable; en fait, il est représentable par le
même schéma $X$, considéré maintenant comme objet de $(\Sch'/S)_{fppf}$, et
\item si $F$ est un espace algébrique sur $S$, alors $f^{-1}F$ est lui aussi un
espace algébrique sur $S$.
\end{enumerate}
\end{lemma}

\begin{proof}
Soit $X \in \Ob((\Sch/S)_{fppf})$. Notons $h_X$ le
faisceau représentable sur $(\Sch/S)_{fppf}$ associé à $X$, et
$h'_X$ le faisceau représentable sur $(\Sch'/S)_{fppf}$ associé à
$X$. D'après la description de $f^{-1}$ dans
Topologies sur les schémas, section \ref{topologies-section-change-alpha},
on a $f^{-1}h_X = h'_X$. Cela prouve (1).

\medskip\noindent
Supposons ensuite que $F$
soit un espace algébrique sur $S$. D'après le Lemme \ref{lemma-space-presentation},
cela signifie que $F = h_U/h_R$ pour une relation d'équivalence \'etale
$R \to U \times_S U$ dans $(\Sch/S)_{fppf}$. Comme $f^{-1}$ est un
foncteur exact, on en conclut que $f^{-1}F = h'_U/h'_R$. Par conséquent,
$f^{-1}F$ est un espace algébrique sur $S$ d'après le Théorème \ref{theorem-presentation}.
\end{proof}

\noindent
Notons que ce lemme est purement ensembliste et n'a pratiquement aucun contenu.
De plus, il n'est pas vrai en général que la restriction d'un espace
algébrique sur le site plus grand soit un espace algébrique sur le site plus petit
(pour de simples raisons de cardinalité). On ne peut donc utiliser un lemme
de ce genre que pour agrandir la catégorie de base, et jamais pour la rétrécir.

\begin{lemma}
\label{lemma-fully-faithful}
Supposons que $\Sch_{fppf}$ soit contenu dans $\Sch'_{fppf}$.
Soit $S$ un objet de $\Sch_{fppf}$. Notons
$\textit{Espaces}/S$ la catégorie des espaces algébriques sur $S$
définie à l'aide de $\Sch_{fppf}$. De même, notons
$\textit{Espaces}'/S$ la catégorie des espaces algébriques sur $S$
définie à l'aide de $\Sch'_{fppf}$. La construction du
Lemme \ref{lemma-change-big-site}
définit un foncteur pleinement fidèle
$$
\textit{Espaces}/S \longrightarrow \textit{Espaces}'/S
$$
dont l'image essentielle est formée des $X' \in \Ob(\textit{Espaces}'/S)$
pour lesquels il existe $U, R \in \Ob((\Sch/S)_{fppf})$\footnote{Exiger
l'existence de $R$ est nécessaire à cause de notre choix de la fonction $Bound$ dans
Théorie des ensembles, Équation (\ref{sets-equation-bound}). La taille du produit fibré
$U \times_{X'} U$ peut croître plus vite que $Bound$ en fonction de la taille de $U$. On
peut l'illustrer en posant $S = \Spec(A)$, $U = \Spec(A[x_i, i \in I])$ et
$R = \coprod_{(\lambda_i) \in A^I} \Spec(A[x_i, y_i]/(x_i - \lambda_i y_i))$.
Dans ce cas, la taille de $R$ croît comme $\kappa^\kappa$, où $\kappa$ est la
taille de $U$.} et des morphismes
$$
U \longrightarrow X'
\quad\text{et}\quad
R \longrightarrow U \times_{X'} U
$$
dans $\Sh((\Sch'/S)_{fppf})$ qui sont surjectifs comme morphismes de faisceaux
(par exemple, si les morphismes affichés sont surjectifs et \'etales).
\end{lemma}

\begin{proof}
Dans Sites et faisceaux, Lemme \ref{sites-lemma-bigger-site}, nous avons vu que le foncteur
$f^{-1} : \Sh((\Sch/S)_{fppf}) \to \Sh((\Sch'/S)_{fppf})$
est pleinement fidèle (voir la discussion dans
Topologies sur les schémas, section \ref{topologies-section-change-alpha}).
Le foncteur affiché dans le lemme est donc pleinement fidèle.

\medskip\noindent
Supposons que $X' \in \Ob(\textit{Espaces}'/S)$ soit tel qu'il existe
$U \in \Ob((\Sch/S)_{fppf})$ et un morphisme $U \to X'$ dans
$\Sh((\Sch'/S)_{fppf})$ qui soit surjectif comme morphisme de faisceaux.
Soit $U' \to X'$ un morphisme \'etale surjectif avec
$U' \in \Ob((\Sch'/S)_{fppf})$. Posons $\kappa = \text{size}(U)$, voir
Théorie des ensembles, section \ref{sets-section-categories-schemes}.
Alors $U$ possède un recouvrement ouvert affine $U = \bigcup_{i \in I} U_i$
avec $|I| \leq \kappa$. Remarquons que $U' \times_{X'} U \to U$ est \'etale
et surjectif. Pour chaque $i$, on peut choisir un ouvert quasi-compact
$U'_i \subset U'$ tel que $U'_i \times_{X'} U_i \to U_i$
soit surjectif (car le schéma $U' \times_{X'} U_i$ est la
réunion des ouverts de Zariski $W \times_{X'} U_i$ pour $W \subset U'$
affine et parce que $U' \times_{X'} U_i \to U_i$ est \'etale, donc ouvert).
Alors $\coprod_{i \in I} U'_i \to X'$ est \'etale et surjectif
en vertu de notre hypothèse selon laquelle $U \to X'$ et donc $\coprod U_i \to X'$
est un morphisme surjectif de faisceaux (détails omis).
Comme $U'_i \times_{X'} U \to U'_i$ est un morphisme surjectif de faisceaux
et que $U'_i$ est quasi-compact,
on peut trouver un ouvert quasi-compact $W_i \subset U'_i \times_{X'} U$
tel que $W_i \to U'_i$ soit surjectif comme morphisme de faisceaux
(détails omis). Alors $W_i \to U$ est \'etale et l'on en conclut
que $\text{size}(W_i) \leq \text{size}(U)$, voir
Théorie des ensembles, Lemme \ref{sets-lemma-bound-finite-type}. D'après
Théorie des ensembles, Lemme \ref{sets-lemma-bound-by-covering}, on en conclut
que $\text{size}(U'_i) \leq \text{size}(U)$.
Ainsi, $\coprod_{i \in I} U'_i$ est isomorphe à un objet de
$(\Sch/S)_{fppf}$ d'après Théorie des ensembles, Lemme \ref{sets-lemma-bound-size}.

\medskip\noindent
Soient maintenant $X'$, $U \to X'$ et $R \to U \times_{X'} U$ comme dans
l'énoncé du lemme. Au paragraphe précédent, nous avons vu que
l'on peut trouver $U' \in \Ob((\Sch/S)_{fppf})$ et un morphisme \'etale surjectif
$U' \to X'$ dans $\Sh((\Sch'/S)_{fppf})$. Alors
$U' \times_{X'} U \to U'$ est un morphisme surjectif de faisceaux, c'est-à-dire
qu'il existe un recouvrement fppf $\{U'_i \to U'\}$ tel que $U'_i \to U'$ se factorise par
$U' \times_{X'} U \to U'$.
D'après Théorie des ensembles, Lemme \ref{sets-lemma-bound-fppf-covering},
on peut trouver
$\tilde U \to U'$
qui soit surjectif, plat et localement de présentation finie,
avec $\text{size}(\tilde U) \leq \text{size}(U')$,
tel que $\tilde U \to U'$ se factorise par
$U' \times_{X'} U \to U'$. Considérons alors
$$
\xymatrix{
U' \times_{X'} U' \ar[d] &
\tilde U \times_{X'} \tilde U  \ar[l] \ar[d] \ar[r] &
U \times_{X'} U \ar[d] \\
U' \times_S U' & \tilde U \times_S \tilde U \ar[l] \ar[r] & U \times_S U
}
$$
Les carrés sont cartésiens. On sait que les objets de la ligne inférieure
sont représentés par des objets de $(\Sch/S)_{fppf}$. D'après le résultat de
l'argument du paragraphe précédent, il en va de même pour $U \times_{X'} U$
(puisque, par hypothèse, on dispose du morphisme surjectif de faisceaux
$R \to U \times_{X'} U$). Comme $(\Sch/S)_{fppf}$ est stable par produits
fibrés (par construction), on voit que $\tilde U \times_{X'} \tilde U$
est représenté par un objet de $(\Sch/S)_{fppf}$. Enfin, le
morphisme $\tilde U \times_{X'} \tilde U \to U' \times_{X'} U'$ est
un morphisme surjectif de faisceaux fppf, puisque $\tilde U \to U'$ l'est.
On peut donc appliquer une fois encore le résultat du paragraphe précédent
pour conclure que $R' = U' \times_{X'} U'$ est représenté par un objet
de $(\Sch/S)_{fppf}$. À ce stade, le
Lemme \ref{lemma-space-presentation} et le
Théorème \ref{theorem-presentation} impliquent que $X = h_{U'}/h_{R'}$
est un objet de $\textit{Espaces}/S$ tel que $f^{-1}X \cong X'$,
comme souhaité.
\end{proof}



\section{Changement de schéma de base}
\label{section-change-base-scheme}

\noindent
Dans cette section, nous examinons brièvement ce qui se passe lorsque l'on
change de schéma de base. Le point essentiel est que, étant donné un morphisme
$S \to S'$ de schémas de base, tout espace algébrique au-dessus de $S$ peut être
considéré comme un espace algébrique au-dessus de $S'$. De plus, étant donné un
espace algébrique $F'$ au-dessus de $S'$, il existe un changement de base $F'_S$, qui est
un espace algébrique au-dessus de $S$. Nous expliquons seulement le cas où $S \to S'$
est un morphisme du grand site fppf considéré. Si seul l'un des schémas $S$ et $S'$
appartient au grand site, on commence par agrandir celui-ci comme dans la
section \ref{section-change-big-site}.

\begin{lemma}
\label{lemma-change-base-scheme}
Soit $\Sch_{fppf}$ un grand site.
Soit $g : S \to S'$ un morphisme de $\Sch_{fppf}$.
Soit $j : (\Sch/S)_{fppf} \to (\Sch/S')_{fppf}$ le
foncteur de localisation correspondant.
Soit $F$ un faisceau d'ensembles sur $(\Sch/S)_{fppf}$.
Alors
\begin{enumerate}
\item pour un schéma $T'$ au-dessus de $S'$, on a
$j_!F(T'/S') =
\coprod\nolimits_{\varphi : T' \to S} F(T' \xrightarrow{\varphi} S),$
\item si $F$ est représentable par un schéma
$X \in \Ob((\Sch/S)_{fppf})$,
alors $j_!F$ est représentable par $j(X)$, c'est-à-dire par
$X$ considéré comme un schéma au-dessus de $S'$, et
\item si $F$ est un espace algébrique au-dessus de $S$, alors $j_!F$ est un
espace algébrique au-dessus de $S'$ et, si $F = U/R$ est une présentation,
alors $j_!F = j(U)/j(R)$ est une présentation.
\end{enumerate}
Soit $F'$ un faisceau d'ensembles sur $(\Sch/S')_{fppf}$. Alors
\begin{enumerate}
\item[(4)] pour un schéma $T$ au-dessus de $S$, on a $j^{-1}F'(T/S) = F'(T/S')$,
\item[(5)] si $F'$ est représentable par un schéma
$X' \in \Ob((\Sch/S')_{fppf})$, alors
$j^{-1}F'$ est représentable, à savoir par $X'_S = S \times_{S'} X'$, et
\item[(6)] si $F'$ est un espace algébrique, alors
$j^{-1}F'$ est un espace algébrique et, si $F' = U'/R'$ est une présentation,
alors $j^{-1}F' = U'_S/R'_S$ est une présentation.
\end{enumerate}
\end{lemma}

\begin{proof}
Les foncteurs $j_!$, $j_*$ et $j^{-1}$ sont définis dans
Sites et faisceaux, Lemme \ref{sites-lemma-relocalize},
où l'on montre également que $j = j_{S/S'}$ est la localisation
de $(\Sch/S')_{fppf}$ en l'objet $S/S'$. Tout ce qui concerne
les foncteurs de localisation s'applique donc à $j$.
La formule de (1) est celle de
Sites et faisceaux, Lemme \ref{sites-lemma-describe-j-shriek-good-site}.
Par définition, $j_!$ est adjoint à gauche du foncteur de restriction $j^{-1}$ ;
$j_!$ est donc exact à droite. D'après
Sites et faisceaux, Lemme \ref{sites-lemma-j-shriek-commutes-equalizers-fibre-products},
il commute en outre aux produits fibrés et aux égalisateurs.
D'après
Sites et faisceaux, Lemme \ref{sites-lemma-describe-j-shriek-representable},
on a $j_!h_X = h_{j(X)}$, d'où (2).
Si $F$ est un espace algébrique au-dessus de $S$, écrivons $F = U/R$
(Lemme \ref{lemma-space-presentation}) ;
on obtient
$$
j_!F = j(U)/j(R)
$$
puisque $j_!$, étant exact à droite, commute aux conoyaux ; de plus,
$j(R) = j(U) \times_{j_!F} j(U)$ puisque $j_!$ commute aux produits fibrés.
Comme les morphismes $j(s), j(t) : j(R) \to j(U)$ ne sont autres que les morphismes
$s, t : R \to U$ (considérés comme des morphismes de schémas au-dessus de $S'$), ils
sont encore \'etales. Ainsi, $(j(U), j(R), s, t)$ est une relation d'équivalence \'etale.
Le
Théorème \ref{theorem-presentation}
permet donc de conclure que $j_!F$ est un espace algébrique.

\medskip\noindent
Démonstration de (4), (5) et (6). La description de $j^{-1}$ se trouve dans
Sites et faisceaux, section \ref{sites-section-localize}.
La restriction du faisceau représentable associé à $X'/S'$
est le faisceau représentable associé à
$X'_S = S \times_{S'} X'$ d'après
Sites et faisceaux, Lemme \ref{sites-lemma-localize-given-products}.
Le foncteur de restriction $j^{-1}$ est exact, donc $j^{-1}F' = U'_S/R'_S$.
Par exactitude encore, le faisceau $R'_S$ demeure une relation d'équivalence sur
$U'_S$. Enfin, les deux morphismes $R'_S \to U'_S$ sont \'etales, car ils se
déduisent par changement de base des morphismes \'etales $R' \to U'$. Ainsi,
$j^{-1}F' = U'_S/R'_S$ est un espace algébrique d'après le
Théorème \ref{theorem-presentation},
ce qui achève la démonstration.
\end{proof}

\noindent
Notons que la présentation $j_!F = j(U)/j(R)$ n'est autre que la présentation
de $F$, considérée comme une présentation par des schémas au-dessus de $S'$.
La définition suivante a donc un sens.

\begin{definition}
\label{definition-base-change}
Soit $\Sch_{fppf}$ un grand site fppf.
Soit $S \to S'$ un morphisme de ce site.
\begin{enumerate}
\item Si $F'$ est un espace algébrique au-dessus de $S'$, alors le
{\it changement de base de $F'$ à $S$} est
l'espace algébrique $j^{-1}F'$ décrit dans le
Lemme \ref{lemma-change-base-scheme}. Nous le notons $F'_S$.
\item Si $F$ est un espace algébrique au-dessus de $S$, alors $F$,
{\it considéré comme un espace algébrique au-dessus de $S'$},
est l'espace algébrique $j_!F$ au-dessus de $S'$ décrit dans le
Lemme \ref{lemma-change-base-scheme}. Nous le notons souvent simplement
$F$ ; sinon, nous écrivons $j_!F$.
\end{enumerate}
\end{definition}

\noindent
L'espace algébrique $j_!F$ est muni d'un morphisme canonique
$j_!F \to S$ d'espaces algébriques au-dessus de $S'$. Cela vient simplement
de ce que le faisceau $j_!F$ admet un morphisme vers $h_S$ (voir, par exemple, la
description explicite du
Lemme \ref{lemma-change-base-scheme}).
En effet, dans
Sites et faisceaux, Lemme \ref{sites-lemma-essential-image-j-shriek},
nous avons vu que la catégorie des faisceaux sur $(\Sch/S)_{fppf}$
est équivalente à la catégorie des couples $(\mathcal{F}', \mathcal{F}' \to h_S)$
formés d'un faisceau sur $(\Sch/S')_{fppf}$ et d'un morphisme de faisceaux
$\mathcal{F}' \to h_S$. L'équivalence associe au faisceau $\mathcal{F}$
le couple $(j_!\mathcal{F}, j_!\mathcal{F} \to h_S)$.
Ceci, joint à ce qui précède, conduit au résultat suivant
pour les catégories d'espaces algébriques.

\begin{lemma}
\label{lemma-category-of-spaces-over-smaller-base-scheme}
Soit $\Sch_{fppf}$ un grand site fppf.
Soit $S \to S'$ un morphisme de ce site.
La construction ci-dessus donne une équivalence de
catégories
$$
\begin{gathered}
\left\{
\begin{matrix}
\text{catégorie des espaces}\\
\text{algébriques sur }S
\end{matrix}
\right\}\\[4pt]
\leftrightarrow\\[4pt]
\left\{
\begin{matrix}
\text{catégorie des couples }(F', F' \to S)\text{ formés}\\
\text{d'un espace algébrique }F'\text{ sur }S'\\
\text{et d'un morphisme }F' \to S\text{ d'espaces algébriques}\\
\text{sur }S'
\end{matrix}
\right\}
\end{gathered}
$$
\end{lemma}

\begin{proof}
Soit $F$ un espace algébrique au-dessus de $S$. Le foncteur de gauche à droite
associe le couple $(j_!F, j_!F \to S)$ à $F$,
et ce couple appartient au membre de droite d'après le
Lemme \ref{lemma-change-base-scheme}.
Puisque ceci définit une équivalence de catégories de faisceaux d'après
Sites et faisceaux, Lemme \ref{sites-lemma-essential-image-j-shriek},
il suffit, pour achever la démonstration, de montrer ceci :
si $F$ est un faisceau et $j_!F$ un espace algébrique, alors $F$
est un espace algébrique. Pour cela, écrivons
$j_!F = U'/R'$ comme dans le
Lemme \ref{lemma-space-presentation},
avec $U', R' \in \Ob((\Sch/S')_{fppf})$.
Les composés $U' \to j_!F \to S$ et $R' \to j_!F \to S$
sont alors des morphismes de schémas au-dessus de $S'$. Notons $U, R$ les
objets correspondants de $(\Sch/S)_{fppf}$. Les deux morphismes
$R' \to U'$ sont des morphismes au-dessus de $S$ et correspondent donc à des
morphismes $R \to U$. Comme ce sont les mêmes
morphismes (considérés au-dessus de $S$), nous obtenons une relation
d'équivalence \'etale au-dessus de $S$. Puisque $j_!$ définit une équivalence de
catégories de faisceaux (voir la référence ci-dessus), on a
$F = U/R$ ; le
Théorème \ref{theorem-presentation}
montre alors que $F$ est un espace algébrique.
\end{proof}

\noindent
Le lemme suivant est une légère reformulation de ce qui précède.

\begin{lemma}
\label{lemma-rephrase}
Soit $\Sch_{fppf}$ un grand site fppf.
Soit $S \to S'$ un morphisme de ce site.
Soit $F'$ un faisceau sur $(\Sch/S')_{fppf}$.
Les assertions suivantes sont équivalentes :
\begin{enumerate}
\item la restriction $F'|_{(\Sch/S)_{fppf}}$
est un espace algébrique au-dessus de $S$, et
\item le faisceau $h_S \times F'$ est un espace algébrique au-dessus de $S'$.
\end{enumerate}
\end{lemma}

\begin{proof}
La restriction et le produit se correspondent par
l'équivalence de catégories de
Sites et faisceaux, Lemme \ref{sites-lemma-essential-image-j-shriek},
de sorte que le
Lemme \ref{lemma-category-of-spaces-over-smaller-base-scheme}
ci-dessus donne le résultat.
\end{proof}

\noindent
Nous terminons cette section par un lemme de compatibilité.

\begin{lemma}
\label{lemma-viewed-as-properties}
Soit $\Sch_{fppf}$ un grand site fppf.
Soit $S \to S'$ un morphisme de ce site.
Soit $F$ un espace algébrique au-dessus de $S$.
Soit $T$ un schéma au-dessus de $S$ et soit $f : T \to F$
un morphisme au-dessus de $S$.
Soit $f' : T' \to F'$ le morphisme au-dessus de $S'$ obtenu à partir de
$f$ en appliquant l'équivalence de catégories décrite dans le
Lemme \ref{lemma-category-of-spaces-over-smaller-base-scheme}.
Pour toute propriété $\mathcal{P}$ comme dans la
Définition \ref{definition-relative-representable-property},
on a $\mathcal{P}(f') \Leftrightarrow \mathcal{P}(f)$.
\end{lemma}

\begin{proof}
Supposons que $U$ soit un schéma au-dessus de $S$ et que $U \to F$ soit un
morphisme \'etale surjectif. Notons $U'$ le schéma $U$ considéré comme un schéma au-dessus de $S'$. Dans le
Lemme \ref{lemma-change-base-scheme},
nous avons vu que $U' \to F'$ est \'etale surjectif. Puisque
$$
j(T \times_{f, F} U) = T' \times_{f', F'} U'
$$
le morphisme de schémas $T \times_{f, F} U \to U$
s'identifie au morphisme de schémas
$T' \times_{f', F'} U' \to U'$.
Il s'agit du même morphisme, considéré au-dessus de schémas de base différents.
Le lemme résulte donc du
Lemme \ref{lemma-representable-morphisms-spaces-property}.
\end{proof}










\begin{multicols}{2}[\section{Autres chapitres}]
\noindent
Préliminaires
\begin{enumerate}
\item \hyperref[introduction-section-phantom]{Introduction}
\item \hyperref[conventions-section-phantom]{Conventions}
\item \hyperref[sets-section-phantom]{Théorie des ensembles}
\item \hyperref[categories-section-phantom]{Catégories}
\item \hyperref[topology-section-phantom]{Topologie}
\item \hyperref[sheaves-section-phantom]{Faisceaux sur les espaces}
\item \hyperref[sites-section-phantom]{Sites et faisceaux}
\item \hyperref[stacks-section-phantom]{Champs}
\item \hyperref[fields-section-phantom]{Corps}
\item \hyperref[algebra-section-phantom]{Algèbre commutative}
\item \hyperref[brauer-section-phantom]{Groupes de Brauer}
\item \hyperref[homology-section-phantom]{Algèbre homologique}
\item \hyperref[derived-section-phantom]{Catégories dérivées}
\item \hyperref[simplicial-section-phantom]{Méthodes simpliciales}
\item \hyperref[more-algebra-section-phantom]{Compléments d'algèbre}
\item \hyperref[smoothing-section-phantom]{Lissage des morphismes d'anneaux}
\item \hyperref[modules-section-phantom]{Faisceaux de modules}
\item \hyperref[sites-modules-section-phantom]{Modules sur les sites}
\item \hyperref[injectives-section-phantom]{Objets injectifs}
\item \hyperref[cohomology-section-phantom]{Cohomologie des faisceaux}
\item \hyperref[sites-cohomology-section-phantom]{Cohomologie sur les sites}
\item \hyperref[dga-section-phantom]{Algèbre différentielle graduée}
\item \hyperref[dpa-section-phantom]{Algèbre à puissances divisées}
\item \hyperref[sdga-section-phantom]{Faisceaux différentiels gradués}
\item \hyperref[hypercovering-section-phantom]{Hyperrecouvrements}
\end{enumerate}
Schémas
\begin{enumerate}
\setcounter{enumi}{25}
\item \hyperref[schemes-section-phantom]{Schémas}
\item \hyperref[constructions-section-phantom]{Constructions de schémas}
\item \hyperref[properties-section-phantom]{Propriétés des schémas}
\item \hyperref[morphisms-section-phantom]{Morphismes de schémas}
\item \hyperref[coherent-section-phantom]{Cohomologie des schémas}
\item \hyperref[divisors-section-phantom]{Diviseurs}
\item \hyperref[limits-section-phantom]{Limites de schémas}
\item \hyperref[varieties-section-phantom]{Variétés}
\item \hyperref[topologies-section-phantom]{Topologies sur les schémas}
\item \hyperref[descent-section-phantom]{Descente}
\item \hyperref[perfect-section-phantom]{Catégories dérivées de schémas}
\item \hyperref[more-morphisms-section-phantom]{Compléments sur les morphismes}
\item \hyperref[flat-section-phantom]{Compléments sur la platitude}
\item \hyperref[groupoids-section-phantom]{Schémas en groupoïdes}
\item \hyperref[more-groupoids-section-phantom]{Compléments sur les schémas en groupoïdes}
\item \hyperref[etale-section-phantom]{Morphismes étales de schémas}
\end{enumerate}
Thèmes de théorie des schémas
\begin{enumerate}
\setcounter{enumi}{41}
\item \hyperref[chow-section-phantom]{Homologie de Chow}
\item \hyperref[intersection-section-phantom]{Théorie de l'intersection}
\item \hyperref[pic-section-phantom]{Schémas de Picard des courbes}
\item \hyperref[weil-section-phantom]{Théories cohomologiques de Weil}
\item \hyperref[adequate-section-phantom]{Modules adéquats}
\item \hyperref[dualizing-section-phantom]{Complexes dualisants}
\item \hyperref[duality-section-phantom]{Dualité pour les schémas}
\item \hyperref[discriminant-section-phantom]{Discriminants et différentes}
\item \hyperref[derham-section-phantom]{Cohomologie de de Rham}
\item \hyperref[local-cohomology-section-phantom]{Cohomologie locale}
\item \hyperref[algebraization-section-phantom]{Géométrie algébrique et formelle}
\item \hyperref[curves-section-phantom]{Courbes algébriques}
\item \hyperref[resolve-section-phantom]{Résolution des surfaces}
\item \hyperref[models-section-phantom]{Réduction semi-stable}
\item \hyperref[functors-section-phantom]{Foncteurs et morphismes}
\item \hyperref[equiv-section-phantom]{Catégories dérivées de variétés}
\item \hyperref[pione-section-phantom]{Groupes fondamentaux de schémas}
\item \hyperref[etale-cohomology-section-phantom]{Cohomologie étale}
\item \hyperref[crystalline-section-phantom]{Cohomologie cristalline}
\item \hyperref[proetale-section-phantom]{Cohomologie pro-étale}
\item \hyperref[relative-cycles-section-phantom]{Cycles relatifs}
\item \hyperref[more-etale-section-phantom]{Compléments de cohomologie étale}
\item \hyperref[trace-section-phantom]{La formule des traces}
\end{enumerate}
Espaces algébriques
\begin{enumerate}
\setcounter{enumi}{64}
\item \hyperref[spaces-section-phantom]{Espaces algébriques}
\item \hyperref[spaces-properties-section-phantom]{Propriétés des espaces algébriques}
\item \hyperref[spaces-morphisms-section-phantom]{Morphismes d'espaces algébriques}
\item \hyperref[decent-spaces-section-phantom]{Espaces algébriques décents}
\item \hyperref[spaces-cohomology-section-phantom]{Cohomologie des espaces algébriques}
\item \hyperref[spaces-limits-section-phantom]{Limites d'espaces algébriques}
\item \hyperref[spaces-divisors-section-phantom]{Diviseurs sur les espaces algébriques}
\item \hyperref[spaces-over-fields-section-phantom]{Espaces algébriques sur des corps}
\item \hyperref[spaces-topologies-section-phantom]{Topologies sur les espaces algébriques}
\item \hyperref[spaces-descent-section-phantom]{Descente et espaces algébriques}
\item \hyperref[spaces-perfect-section-phantom]{Catégories dérivées d'espaces}
\item \hyperref[spaces-more-morphisms-section-phantom]{Compléments sur les morphismes d'espaces}
\item \hyperref[spaces-flat-section-phantom]{Platitude sur les espaces algébriques}
\item \hyperref[spaces-groupoids-section-phantom]{Groupoïdes dans les espaces algébriques}
\item \hyperref[spaces-more-groupoids-section-phantom]{Compléments sur les groupoïdes dans les espaces}
\item \hyperref[bootstrap-section-phantom]{Amorçage}
\item \hyperref[spaces-pushouts-section-phantom]{Sommes amalgamées d'espaces algébriques}
\end{enumerate}
Thèmes de géométrie
\begin{enumerate}
\setcounter{enumi}{81}
\item \hyperref[spaces-chow-section-phantom]{Groupes de Chow des espaces}
\item \hyperref[groupoids-quotients-section-phantom]{Quotients de groupoïdes}
\item \hyperref[spaces-more-cohomology-section-phantom]{Compléments sur la cohomologie des espaces}
\item \hyperref[spaces-simplicial-section-phantom]{Espaces simpliciaux}
\item \hyperref[spaces-duality-section-phantom]{Dualité pour les espaces}
\item \hyperref[formal-spaces-section-phantom]{Espaces algébriques formels}
\item \hyperref[restricted-section-phantom]{Algébrisation des espaces formels}
\item \hyperref[spaces-resolve-section-phantom]{Résolution des surfaces revisitée}
\end{enumerate}
Théorie des déformations
\begin{enumerate}
\setcounter{enumi}{89}
\item \hyperref[formal-defos-section-phantom]{Théorie formelle des déformations}
\item \hyperref[defos-section-phantom]{Théorie des déformations}
\item \hyperref[cotangent-section-phantom]{Le complexe cotangent}
\item \hyperref[examples-defos-section-phantom]{Problèmes de déformation}
\end{enumerate}
Champs algébriques
\begin{enumerate}
\setcounter{enumi}{93}
\item \hyperref[algebraic-section-phantom]{Champs algébriques}
\item \hyperref[examples-stacks-section-phantom]{Exemples de champs}
\item \hyperref[stacks-sheaves-section-phantom]{Faisceaux sur les champs algébriques}
\item \hyperref[criteria-section-phantom]{Critères de représentabilité}
\item \hyperref[artin-section-phantom]{Axiomes d'Artin}
\item \hyperref[quot-section-phantom]{Espaces de Quot et de Hilbert}
\item \hyperref[stacks-properties-section-phantom]{Propriétés des champs algébriques}
\item \hyperref[stacks-morphisms-section-phantom]{Morphismes de champs algébriques}
\item \hyperref[stacks-limits-section-phantom]{Limites de champs algébriques}
\item \hyperref[stacks-cohomology-section-phantom]{Cohomologie des champs algébriques}
\item \hyperref[stacks-perfect-section-phantom]{Catégories dérivées de champs}
\item \hyperref[stacks-introduction-section-phantom]{Introduction aux champs algébriques}
\item \hyperref[stacks-more-morphisms-section-phantom]{Compléments sur les morphismes de champs}
\item \hyperref[stacks-geometry-section-phantom]{La géométrie des champs}
\end{enumerate}
Thèmes de théorie des espaces de modules
\begin{enumerate}
\setcounter{enumi}{107}
\item \hyperref[moduli-section-phantom]{Champs de modules}
\item \hyperref[moduli-curves-section-phantom]{Espaces de modules de courbes}
\end{enumerate}
Divers
\begin{enumerate}
\setcounter{enumi}{109}
\item \hyperref[examples-section-phantom]{Exemples}
\item \hyperref[exercises-section-phantom]{Exercices}
\item \hyperref[guide-section-phantom]{Guide bibliographique}
\item \hyperref[desirables-section-phantom]{Desiderata}
\item \hyperref[coding-section-phantom]{Style de programmation}
\item \hyperref[obsolete-section-phantom]{Obsolète}
\item \hyperref[fdl-section-phantom]{Licence de documentation libre GNU}
\item \hyperref[index-section-phantom]{Index généré automatiquement}
\end{enumerate}
\end{multicols}


\bibliography{my}
\bibliographystyle{amsalpha}

\end{document}
